\documentclass[11pt,reqno]{amsart}
\usepackage[T1]{fontenc}
\usepackage{lmodern,amsmath,amssymb,mathtools,microtype,booktabs,array,longtable}
\usepackage[margin=1in]{geometry}
\usepackage{graphicx}
\usepackage{xr-hyper}
\usepackage[hidelinks]{hyperref}
\ifdefined\pdfinfoomitdate\pdfinfoomitdate=1\fi
\ifdefined\pdftrailerid\pdftrailerid{}\fi
\ifdefined\pdfsuppressptexinfo\pdfsuppressptexinfo=15\fi
\numberwithin{equation}{section}
\newtheorem{theorem}{Theorem}[section]
\newtheorem{proposition}[theorem]{Proposition}
\newtheorem{lemma}[theorem]{Lemma}
\newtheorem{corollary}[theorem]{Corollary}
\theoremstyle{definition}\newtheorem{definition}[theorem]{Definition}
\newtheorem{example}[theorem]{Example}
\theoremstyle{remark}\newtheorem{remark}[theorem]{Remark}
\newcommand{\E}{\mathbb E}\newcommand{\Prb}{\mathbb P}
\newcommand{\R}{\mathbb R}\newcommand{\one}{\mathbf1}
\newcommand{\Q}{\mathcal Q}\newcommand{\TV}{\mathrm{TV}}
\newcommand{\cp}{\mathrm{cp}}\newcommand{\ev}{\mathrm{ev}}\newcommand{\st}{\mathrm{st}}
\DeclareMathOperator{\tr}{tr}\DeclareMathOperator{\dist}{dist}
\DeclareMathOperator{\law}{Law}\DeclareMathOperator{\rank}{rank}
\DeclareMathOperator{\supp}{supp}
\emergencystretch=2em

\externaldocument[][nocite]{supplement-refs}[General_Theta_Foundations_I_R66_Supplement.pdf]
\hypersetup{pdftitle={General Theta Foundations I: Acquired Geometry and Causal Resource Transfer},pdfauthor={Qian Qi}}
\title[Acquired geometry and causal resource transfer]{General Theta Foundations I:\\Acquired Geometry and Causal Resource Transfer}
\author{Qian Qi}\date{October 2026}
\subjclass[2020]{Primary 60G40, 62C20; Secondary 68Q45, 81P15, 90C22, 94A29}
\keywords{causal experiment, performed stopped law, boundary optimality, finite memory, decision quotient, quantum instrument, robust prediction}
\begin{document}
\begin{abstract}
We study exact finite classical memory at the least feasible cost in prepared causal experiments. The general performed-law theory gives full-history boundary attainment but does not imply finite-state realization. We derive a weighted minimum-cost potential identity for single-command quantum or classical instruments with state- and model-dependent reports. A support condition on cost-maximizing models forces first-herald commitment. On this boundary, compact minimax identifies a probability-valued optimal decision face without assuming unique Bayes actions. Exact stochastic memory is equivalent to an anchored finite convex family invariant under report shifts; its smallest generator count is the least live memory. This replaces the deterministic residual quotient in the robust tied case. An unknown-random-horizon calibration task forces replication of every Born probability by one randomized response law, even though all wordwise Bayes actions are tied. Protected classical channels realize every prescribed finite exact memory dimension. For an algebraic rotating qubit, complete depolarization needs two labels, every strictly damped orbit has a finite exact stochastic realization, and the undamped orbit has none despite Hankel rank three. We determine the exact global three-label damping threshold and give a terminating algebraic procedure for minimal exact memory in the damped family. The previously proved unique-decision obstruction, two-label squared-prediction frontier, performed-law duality, timing and actual-history geometric transfer remain intact with their separate hypotheses. Finite linear rank, positive stochastic realization, protected physical dimension and observer memory are distinct resources.
\end{abstract}
\maketitle
\section{Introduction and current theorem structure}\label{sec:current-overview}\label{sec:introduction}
There are three different questions at a finite expected-cost boundary: whether a full-history optimum is attained, whether finite-memory policies approximate it at the same budget, and whether an exact finite-memory optimum exists. The performed-law compactness theorem answers the first on a general finite-branching interface. The heralded results answer the latter two on more structured physical interfaces. None of these statements identifies a finite physical Hilbert dimension with a finite classical controller.

Our principal results concern exact stochastic boundary implementation. Theorem~\ref{thm:weighted-rigidity} allows state- and model-dependent reports and derives first-herald rigidity from an actual cost potential plus a declared support condition. Theorem~\ref{thm:weighted-face} identifies the full robust optimal face, retaining modelwise inequalities when Bayes actions tie. Theorem~\ref{thm:convex-memory} characterizes its finite stochastic implementation by an anchored shift-invariant convex family. This is a positive-realization criterion, not a newly invented deterministic residual theorem. It recovers the earlier unique-decision result while permitting a genuinely randomized robust optimum.

The calibration application makes the distinction concrete. An unknown geometric observation horizon and two opposed calibration losses force one common probability response equal to the protected system's Born probabilities at every failure count. Theorem~\ref{thm:calibration-k} gives exact memory $k$ for every $k\ge2$. Theorem~\ref{thm:qubit-positive-phase} gives a sharp global three-label threshold for the algebraic qubit, finite exact realization at every strict damping, and impossibility at zero damping despite finite Hankel rank. Proposition~\ref{prop:finite-realization-test} turns the finite size bound into minimal algebraic synthesis. This is a different task from the preserved squared-sign prediction problem, whose general intermediate-memory curve remains open.

The exact-memory machinery has classical antecedents in residual automata, linear stochastic representations and positive realization \cite{RabinScott59,Jaeger00,ThonJaeger15,MonrasWinter14}. Our assertions are the cost-face and optimal-face derivation, their common-model implementation quantifiers, and the exact memory conclusions for the declared physical tasks. These theorems do not classify arbitrary POMDP controllers or certify earliest priority for their ingredients.

The performed-law, primal--dual, timing and actual-history transfer proofs are included below; their full technical consequences are given in the supplement. These complementary results have separate hypotheses and are not inferred from the calibration example.


The sparse witness count in the performed-law dual is not an implementation-memory count. Absorbing-MDP extreme occupation measures and finite mixtures of deterministic stationary policies are primal observed-state objects \cite{PiunovskiyZhang24,DufourMixtures}; our cost/risk preparations are dual physical witnesses. A stationary policy can read a continuous state, while the observer here must act from counted retained labels without that input. The general geometric transfer still uses its actual-history lower witness and one common paid upper program. Neither it nor the separate timing curve is replaced by the calibration example.

\section{Prepared kernels, acquired task laws and observers}\label{sec:data}
All hidden-state, action and report spaces are standard Borel. A fixed unknown index $\theta\in\Theta$ specifies an actual preparation and positive normalized causal kernels. Positivity means nonnegative measures, not an everywhere strictly positive density. Dependence on the entire physical past may be included in the hidden state. A legal policy sees past visible reports and its admitted randomization, but not $\theta$ or a future audit. Iteration of these kernels defines its finite-prefix laws; the consistent countable path law is their kernel extension.

Each call has a cost $c_i\ge c_*>0$. The decision to make the next call is predictable. In the geometric route, the declared productive-call flag is also predictable; in the primitive-instrument route, productive reports are specified outcomes and can be random. Commitment is permitted only after a productive report, and occurs exactly once almost surely in each model. At that time a specified future audit $Z\in H$ is scored, where $H\subset\R^k$ is nonempty, compact and convex. The prediction is committed before that audit is observed. A prospective joint report--audit kernel is defined at every productive call, including rejected calls. The stopping coin is independent of the audit conditional on the entire current visible report and the past. Write $L=\operatorname{diam}(H)^2$; the case $L=0$ is trivial whenever a legal rule exists.

The core cost $C_\tau$ includes initial preparation charged to the core, failed calls and all productive calls through commitment. A separately executed final audit must be added to the raw budget. This convention is fixed before taking a risk infimum. Physical time, stored program constants and temporary computation are not identified with call count.

\subsection*{Three resource interfaces}
A checkpoint procedure may use its whole precommit history, but must output one element of a fixed dictionary of at most $M$ values. This counts its output width only; it is not a claim of $\log_2 M$ total working bits.

An autonomous observer has an alphabet $S$, $|S|\le M$, and one fixed parameter-independent program. At an intercall boundary its only persistent observer memory is its label $s\in S$. Its next command distribution is a fixed function of $s$. A physical port name, elapsed time, subprogram index, random seed or former report is not an extra free input. During a call its current command, current report and fresh coins are available. Whatever persists to another call must be encoded in $S$. A fixed decoder $r:S\to H$ reads only the committed label.

In the \emph{event-stopped} interface a post-report kernel emits $(s',e)\in S\times\{0,1\}$. An event $e=1$ ends the experiment and retains $s'$ for the decoder; no subsequent command is executed. This public commitment event is not a stored bit used in continuing computation. The same label may be used live on one path and terminal on another. In the \emph{state-stopped} interface a fixed map $h_*:S\to\{0,1\}$ determines halt status. The initial label is live, and type-one labels are globally absorbing terminal labels with no executed command rows. Entry into a terminal label is allowed only after a productive report. Thus live and terminal types are disjoint in this interface, and only in this interface by definition.

Write $R_i^\Theta(N,M)$, $i\in\{\cp,\ev,\st\}$, for the infimum of worst-case expected squared loss over programs with robust almost-sure commitment and $\sup_\theta\E_\theta C_\tau\le N$. An empty legal class has value $+\infty$. A program chosen once at random is not a free persistent seed. The ideal measurable class may have fixed real constants; effective implementation is stated separately in Section~\ref{sec:composition}.

\subsection*{A physically acquired comparison law}
Execute a declared prepared comparison experiment and obtain its joint report--audit probability $K_0(dx,dz)$ on $X\times H$. Its report marginal is $\nu$. The law is genuinely preparable and its specification has its own resource account. It need not be a member of $\Theta$. In particular no finite set of calibration data is silently promoted to a uniform certificate. Set
\begin{equation}\label{eq:derived}
 f(x)=\int zK_0(dz\mid x),\quad
 w(x)=\int\|z-f(x)\|^2K_0(dz\mid x),\quad
 d_r(x)=w(x)+\|f(x)-r\|^2.
\end{equation}
Disintegration gives Borel versions with $f\in H$; null sets may be completed arbitrarily. The profile is
\begin{equation}\label{eq:profile}
 \Q_m(b)=\min_{\substack{C\subset H,\ 1\le|C|\le m\\
                  0\le h\le b,\ \int h\,d\nu=1}}
        \int [w+\dist(f,C)^2]h\,d\nu,
 \qquad b\ge1,\ m\ge1.
\end{equation}
We use $+\infty$ for an infeasible cap or width. The quantities $f,w$ are task coordinates of acquired predictive information, not a complete substitute for the predictive quotient. Their derivation and the quotient construction in Appendix~\ref{sec:quotient} precede any use of local dimension.

\subsection*{Lower certificates on actual histories}\label{sec:lower-cert}
Fix \emph{one} actual model $\theta_*\in\Theta$, independent of the competing program. At each actual reached productive pre-call history $h$ and legal command $a$, write $P_{h,a}(dy,dz)$ for its conditional full current report--audit law. Let $P^0_{h,a}$ be a measurable comparison row on the same report space, with a Borel mark $x=\pi_{h,a}(y)$ of law $\nu$ and
\begin{equation}\label{eq:full-moments}
 \E^0[Z\mid y]=f(\pi_{h,a}y),\qquad
 \E^0[\|Z\|^2\mid y]=w(\pi_{h,a}y)+\|f(\pi_{h,a}y)\|^2.
\end{equation}
These are conditional identities given the \emph{full} report, not just its mark. Versions and inequalities are required on actual reachable histories, not only almost everywhere under an unrelated comparison path law.

Fix $0<l_-\le1\le u_-<\infty$. For all Borel $q:Y\to[0,1]$ and all finite-image Borel $r:Y\to H$, suppose predictable nonnegative defects $a_-(h,a),e_-(h,a)$ obey
\begin{align}
 \int q\,d(P_{h,a})_Y
 &\le u_-\int q\,d(P^0_{h,a})_Y+a_-(h,a),\label{eq:lower-mass}\\
 \int q(y)\|z-r(y)\|^2P_{h,a}(dy,dz)
 &\ge l_-\int q(y)\|z-r(y)\|^2P^0_{h,a}(dy,dz)-e_-(h,a).
 \label{eq:lower-task}
\end{align}
Independent current randomization can be adjoined to both rows: integrating inequalities for each fixed random value proves the same inequalities with the same defects. All command-conditioned lower inequalities hold before averaging the chosen command.

Let $A_\tau$ count productive calls. At this same actual model require, for every legal competing procedure,
\begin{equation}\label{eq:effort}
 \E_* C_\tau\ge\kappa_-\E_* A_\tau,\qquad\kappa_->0.
\end{equation}
At budget $N$ suppose, uniformly over all robustly feasible competitors,
\begin{equation}\label{eq:occupation}
 \E_*\sum_{j\le\tau:\ j\ {\rm productive}}a_-(H_{j-1},a_j)\le A_N,
 \qquad
 \E_*\sum_{j\le\tau:\ j\ {\rm productive}}e_-(H_{j-1},a_j)\le E_N.
\end{equation}
These are actual occupation certificates, not properties merely of the proposed upper program. Uniform row defects give $A_N\le(N/\kappa_-)\bar a_-$ and $E_N\le(N/\kappa_-)\bar e_-$. Lemma~\ref{lem:drift} supplies a nonuniform alternative.

\subsection*{One common upper construction}\label{sec:upper-cert}
Independently, let $\mathcal G$ be a parameter-independent program which, from each of its start histories, reaches one productive report. Its conditional expected attempt cost is at most $\kappa_+>0$ on every surviving history in every $\theta\in\Theta$. It uses $B$ live roles and returns to its start label on rejection; no reset of the physical state is assumed. It forms a virtual mark $x$, a current choice $J$ and a permitted export $T_Jr$ of its task readout. For every finite dictionary $C$ assume
\begin{equation}\label{eq:exports}
 |\{T_jr:j\in J,\ r\in C\}|\le D_0|C|,
 \qquad D_0\in\mathbb N.
\end{equation}
Every continuing phase and coin needed for this construction is included in the $B$ live roles. For fixed $0<l_+\le1\le u_+<\infty$ and $a_+,e_+\ge0$, assume for all Borel $q:X\to[0,1]$ and finite-image Borel $r:X\to H$, conditionally on every such history,
\begin{align}
 \E[q(x)]&\ge l_+\int q\,d\nu-a_+,\label{eq:upper-mass}\\
 \E[q(x)\|Z-T_Jr(x)\|^2]
 &\le u_+\int q(x)d_{r(x)}(x)\nu(dx)+e_+.
 \label{eq:upper-task}
\end{align}
The expectation here includes the within-attempt choice $J$. Unlike the lower, the upper is not required to hold for each choice separately. The choices retained after this call are nevertheless charged through the actual export or live role. These hypotheses express one common causal construction, not modelwise optimizers.

\section{A constructive finite-budget frontier from primitive instruments}\label{sec:primitive-risk}
The geometric transfer theorem identifies a geometric curve when its converse and construction certificates are available. We now give a different, intrinsic route for finite time-homogeneous experiments: derive a proper fallback from the raw instruments and compute converging bounds on the optimal task risk itself. Neither a reset state, a balancing group, a supplied reference member nor a posterior stored by the controller is assumed. The finite and all-initial-state qualifications are essential. The qualitative invariant-subspace obstruction has a direct antecedent in quantum reachability \cite[Theorem 3.6]{Ying14pre}; it is proved below to obtain the quantitative fallback required by the new budget-preserving approximation.

\subsection{Primitive interface and universal termination}
The unknown model belongs to a nonempty finite set $\Theta$. Model $\theta$ has physical space $\mathcal H_\theta$ of dimension $d_\theta$. The nonempty finite command set $\mathcal A$ is globally legal. A command has completely positive report maps $\mathcal I_{\theta,a,y}$ with normalized sum. Reports range over finite sets, with a declared productive subset $Y_a^+$. After a productive report a controller may commit a decision from a fixed dictionary $C$ or reject; after any rejected or nonproductive report the post-report state is the next state. No additional return or reset is available. An interface requiring extra physical return calls must first represent and charge those calls explicitly; the theorem is not applied by hiding them in an unpriced transition. Commitment occurs exactly once, before a future audit is observed. The audit loss for decision $r\in C$ is represented on the post-report state by
\[
 0\preccurlyeq F_{\theta,a,y,r}\preccurlyeq LI,
 \qquad 0<L<\infty.
\]
Nonnegative call costs have observables $c_*I\preccurlyeq G_{\theta,a}\preccurlyeq c^*I$, with $0<c_*\le c^*<\infty$. These are the precommit core costs. A final audit with one fixed separately charged cost is added to every total budget after the theorem; an action-dependent unmodelled terminal charge is not included by assertion. Policies share one set of command and decision rules across models and unknown initial density matrices. Let $R_C(N)$ be their least worst-model, worst-initial-state expected terminal loss at uniform expected cost at most $N$. Their precommit memory is unrestricted for this value; a finite-label implementation is derived, not assumed.

Fix any command distribution $\pi_a>0$. The fallback chooses that row independently at every call and commits at the first productive report, using one fixed $r_0\in C$. Its killed physical map is
\[
 \mathcal T_\theta=\sum_a\pi_a\sum_{y\notin Y_a^+}\mathcal I_{\theta,a,y},
 \quad B_\theta=I-\mathcal T_\theta^*I,
 \quad O_{\theta,n}=I-(\mathcal T_\theta^*)^nI.
\]
A common dark subspace in one model is a nonzero subspace annihilated by every productive effect and invariant under every Kraus operator of every nonproductive report map. This condition is independent of the Kraus representation: it says that every positive state on the subspace remains there on failure and has zero productive probability for every command.

\begin{theorem}[Primitive fallback and its sharp depth test]\label{thm:primitive-fallback}
The following are equivalent: there is one common policy proper at finite uniform expected cost from every model and initial density matrix; no model has a common dark subspace; and $O_{\theta,d_\theta}$ is positive definite for every $\theta$. In that case the displayed one-row fallback is proper. With
\[
 \delta_\theta=\lambda_{\min}(O_{\theta,d_\theta})>0,
 \qquad \kappa=\max_\theta c^*d_\theta/\delta_\theta,
\]
it costs at most $\kappa$ from every current physical state in every model. One may replace this bound by the sharper certified value
\[
 \kappa=\max_\theta\left\|(I-\mathcal T_\theta^*)^{-1}
                  \sum_a\pi_aG_{\theta,a}\right\|.
\]
Only one live label and one output value are needed; these share a label in event stopping and require two labels in state stopping. Depth $d$ cannot be replaced uniformly by $d-1$. There is no dimension-only upper bound on expected cost.
\end{theorem}
\begin{proof}
Work in one model. Put $K_n=\ker O_n$, so $K_0=\mathcal H$ and
$O_{n+1}=B+\mathcal T^*O_n$. For a Kraus representation $\mathcal T(A)=\sum_vL_vAL_v^*$, positivity gives
\[
 K_{n+1}=\ker B\cap\bigcap_vL_v^{-1}K_n.
\]
The spaces decrease. Equality at two consecutive levels makes that space invariant under every $L_v$ and annihilated by $B$, hence makes all later levels equal. There are at most $d$ strict dimension drops, so $K_d$ is stable. Since every $\pi_a$ is positive, $\ker B$ is the intersection of the productive-effect kernels, and the $L_v$ include nonzero multiples of every nonproductive Kraus operator. Thus a nonzero $K_d$ is a common dark subspace. Conversely every such subspace lies in all $K_n$. Starting in it, no history policy can encounter a productive report: induction holds commandwise and reportwise. Such a policy cannot commit legally.

If $K_d=0$, positivity yields $(\mathcal T^*)^dI\preccurlyeq(1-\delta)I$. Iteration bounds survival after $kd$ fallback calls by $(1-\delta)^k$ for every normalized input. Summing blocks gives expected calls at most $d/\delta$ and expected cost at most $c^*d/\delta$. The finite model set makes the maximum bound finite. The same estimate implies convergence of the positive cost series, whose sum is the stated resolvent; its largest eigenvalue is the exact uniform cost for the fixed row. The fallback does not read the model or physical state. Its single row and fixed decoder give the label counts.

For sharp depth, on basis $e_1,\ldots,e_d$ take one nonproductive Kraus operator $L=\sum_{j=2}^d|e_{j-1}\rangle\langle e_j|$ and a productive effect $|e_1\rangle\langle e_1|$. Then $O_n$ is the projection onto the first $\min(n,d)$ basis vectors. The input $e_d$ requires exactly $d$ unit-cost calls. For the lack of a cost bound, already in dimension one a productive probability $\epsilon>0$ with unit costs has expected time $1/\epsilon$. This proves both limitations.
\end{proof}

The qualitative obstruction is inherited mathematical machinery, not a new discovery of invariant subspaces. Its purpose here is to supply an actual common fallback and its paid constant. Absence of an invariant dark subspace on a restricted preparation is a different problem; the all-initial-state hypothesis cannot be silently removed. Finiteness of $\Theta$ likewise cannot be replaced by modelwise termination in an arbitrary family without a uniform bound.

\subsection{Finite problems and restoration of the exact budget}
For a horizon $T\ge1$, use one common realization plan on the action--report tree, as in Lemma~\ref{lem:common-plan}. At every productive report include either commitment to $r\in C$ or continuation. Let $C_{\theta,T}$, $S_{\theta,T}$ and $V_{\theta,T}$ be, respectively, its expected cost through the first $T$ calls, its open effect after the depth-$T$ decision, and its accumulated terminal-loss operator. They are affine functions of the common sequence weights, computed by composing the primitive maps. In particular $0\preccurlyeq S_{\theta,T}\preccurlyeq I$ and $0\preccurlyeq V_{\theta,T}\preccurlyeq L(I-S_{\theta,T})$.

Define the finite semidefinite value
\begin{equation}\label{eq:primitive-finite}
 q_{C,T}(N)=\min\left\{q:\ C_{\theta,T}+c_*S_{\theta,T}\preccurlyeq NI,
 \quad V_{\theta,T}\preccurlyeq qI\quad(\theta\in\Theta)\right\}.
\end{equation}
The minimum is over the common plan and $0\le q\le L$. The reserve $c_*S$ is not a cutoff declaration of success: every open branch owes another paid call. For $N\ge\kappa$, prefixes of the actual fallback show feasibility. Compactness gives attainment. Set
\[
 D_T(N)=\min\{1,(N/c_*-1)/T\},\qquad
 e_T=(\kappa-c_*)D_T(N).
\]

\begin{theorem}[Same-budget constructive risk frontier]\label{thm:primitive-risk}
Under the preceding primitive conditions with no common dark subspace, the finite values are nondecreasing in $T$ and satisfy $q_{C,T}(N)\le R_C(N)$. For every $N>\kappa$,
\begin{equation}\label{eq:primitive-gap}
 0\le R_C(N)-q_{C,T}(N)
 \le L D_T(N)+\frac{L e_T}{N-\kappa+e_T}.
\end{equation}
The upper is realized by a common proper finite-label observer with expected cost at most the original $N$ in every model and from every initial state. Thus $q_{C,T}(N)\uparrow R_C(N)$ with an explicit $O(T^{-1})$ bound at each fixed $N>\kappa$. No reset or scalar posterior dynamics are required.

Writing $a=|\mathcal A|$ and $y$ for the largest report count, a sufficient live-label bound is
\[
 2+\sum_{t=0}^{T-1}(ay)^t.
\]
The event-stopped total is at most the maximum of this bound and $|C|$; the state-stopped total is at most their sum. These are existence bounds, not optimal memory, running-time or program-bit bounds.

At $N=\kappa$ the completed finite plan still gives cost at most $N+e_T$ and loss at most $q_{C,T}(N)+LD_T(N)$. The same-budget claim at that boundary is not inferred from the displayed fraction. If $\kappa=c_*$, no budget excess is present and the same-budget bound $LD_T(N)$ holds for every $N\ge\kappa$.
\end{theorem}
\begin{proof}
A genuine cost-$N$ policy restricted to depth $T$ satisfies the reserved cost inequality, because after every open path it must make at least one further call. Its truncated terminal loss is no greater than its full nonnegative loss. This proves the lower. A feasible plan at $T+1$ restricts to one at $T$: the new calls cost at least $c_*S_T$, new terminations reduce open mass, and their loss contributions are nonnegative. The reserve and risk inequalities therefore survive restriction, proving monotonicity.

For any feasible prefix, normalization and the conditional cost floor give
\[
 C_T\succeq c_*\sum_{t=0}^{T-1}S_t
       \succeq c_*[I+(T-1)S_T],\qquad S_0=I.
\]
Here each difference of successive open effects is positive. With the reserve this gives $S_T\preccurlyeq D_T(N)I$ in every model. Complete an optimizing prefix by executing the actual fallback at every open leaf. The fallback bound holds at every post-report state, including states never reached by a reference execution. Its cost and loss satisfy
\[
 C^{\rm comp}\preccurlyeq C_T+\kappa S_T
       \preccurlyeq(N+e_T)I,\qquad
 V^{\rm comp}\preccurlyeq V_T+LS_T
       \preccurlyeq(q_{C,T}+LD_T)I.
\]
It is proper because the finite prefix is followed by a proper fallback. There is no assertion that the original policy's expected remaining cost is uniformly small.

For $N>\kappa$, choose the completed policy with probability
$p=(N-\kappa)/(N+e_T-\kappa)$ and the fallback with probability $1-p$, storing the branch in the ensuing live role. Its cost is at most $p(N+e_T)+(1-p)\kappa=N$. Its risk is at most
$p(q_{C,T}+LD_T)+(1-p)L\le q_{C,T}+LD_T+Le_T/(N-\kappa+e_T)$.
If $e_T=0$ choose the completed policy directly. This proves the upper and convergence. The boundary statements follow from the unmixed completion.

A prefix stores its current action--report node, at most $\sum_{t<T}(ay)^t$ live roles; its current report is within-call input. One fallback role and one initial mixture role suffice. At that initial call, sample the branch and its command jointly; after the report either commit or enter the selected branch's counted next role. If both branches chose the same command, their conditional branch weights are proportional to their command probabilities, independent of the physical state, since the executed report instrument is identical. Equivalently use the joint within-call coin directly. No branch bit survives outside the retained label. Commitment uses the same fixed dictionary in both branches. Lemma~\ref{lem:roles} gives the two total counts.
\end{proof}

\begin{corollary}[Compact decisions at a fixed output width]\label{cor:primitive-width}
Suppose the decision carrier $D$ is compact metric and the raw audit losses obey the common decision modulus $\omega$ of \eqref{eq:loss-modulus}. Fix width $m$. For a finite $\eta$-net $D_\eta\subset D$, let $q_{m,T,\eta}(N)$ be the least value in \eqref{eq:primitive-finite} among dictionaries in $D_\eta$ of size at most $m$. Then, for $N>\kappa$,
\[
 [q_{m,T,\eta}(N)-\omega(\eta)]_+
 \le R_m(N)\le q_{m,T,\eta}(N)+LD_T(N)
                         +\frac{Le_T}{N-\kappa+e_T}.
\]
The constructed observer uses at most $m$ terminal values. As $T\to\infty$ and $\eta\downarrow0$ this gives a two-parameter constructive characterization of the unrestricted-history finite-output risk, not just a verifier for a supplied balancing program. With algebraic primitive data and polynomial score dependence on algebraic net points, every finite test is a finite real-algebraic problem; exact real samplers and program arithmetic remain separate implementation resources.
\end{corollary}
\begin{proof}
Round each point of any competing dictionary to the same chosen nearest net point at every history. Neither stopping nor cost changes, and at most $m$ values remain. Its expected loss changes by at most $\omega(\eta)$, because terminal events have total mass one. Restrict this rounded program to a finite prefix to obtain the lower. For the upper, choose an attaining net dictionary and apply Theorem~\ref{thm:primitive-risk}; the fallback uses a point of this same dictionary. There are finitely many such dictionaries, so the optimization is finite. With the stated algebraic inputs all coefficients of the finite plan are algebraic, and semidefinite inequalities can be expressed by principal minors. This gives decidability via the already credited real-algebraic method, not an efficient uniform algorithm for arbitrary real input data.
\end{proof}

\subsection{Removing the crude fallback threshold on every interior budget}
The cost of an arbitrary positive command row may greatly exceed the least feasible cost. It is not an intrinsic resource threshold. The following preliminary optimization removes that restriction without assuming the optimal policy is already known.

Let $\kappa_0$ be the uniform cost bound of the one-row fallback and let $H_\theta$ be its exact cost operator, so $0\preccurlyeq H_\theta\preccurlyeq \kappa_0I$. Ignore terminal loss for this preliminary step and commit one fixed dummy decision. Write $N_{\min}$ for the infimum of uniform expected cost among proper common policies. For a finite prefix define
\[
 W_{\theta,T}(H_\theta)
   =\sum_{h:\,|h|=T}\gamma_{\rm open}(h)\,
                       \mathcal I_{\theta,h}^*(H_\theta).
\]
Here $\gamma_{\rm open}$ includes the decision to continue after the depth-$T$ report, and $\mathcal I_{\theta,h}$ is its raw path map. This is the exact cost of executing the fallback at the open leaves. Both $W$ and the earlier $C_T,S_T$ are affine in the one common realization plan. Set
\begin{align*}
 a_T&=\min\{a:C_{\theta,T}+c_*S_{\theta,T}\preccurlyeq aI\quad(\theta\in\Theta)\},\\
 u_T&=\min\{u:C_{\theta,T}+W_{\theta,T}(H_\theta)\preccurlyeq uI\quad(\theta\in\Theta)\}.
\end{align*}

\begin{theorem}[Constructive feasibility threshold and all interior budgets]\label{thm:interior-budget}
Both finite problems attain their values, and
\[
 c_*\le a_T\uparrow N_{\min},\qquad
 \kappa_0\ge u_T\downarrow N_{\min},
\]
with the quantitative enclosure
\begin{equation}\label{eq:cost-enclosure}
 0\le u_T-a_T\le(\kappa_0-c_*)\min\{1,(\kappa_0/c_*-1)/T\}.
\end{equation}
The least common proper cost $N_{\min}$ is attained, though no finite-memory attainment is concluded at that boundary.

For every $N>N_{\min}$ a finite search through these upper problems eventually exhibits $u_{T_0}<N$. Its optimizing prefix completed by the one-row fallback is an actual finite-label fallback of cost $\kappa=u_{T_0}<N$, valid from every initial state and model. Using it in Theorem~\ref{thm:primitive-risk} gives the same-budget risk enclosure for \emph{every strictly interior feasible budget}, not only budgets above the original one-row cost. If this improved fallback has $B$ live labels, the sufficient live bound becomes
\[
 1+B+\sum_{t=0}^{T-1}(ay)^t.
\]
All decisions may still use one prescribed dictionary; the cost search is independent of its values. Algebraic primitive data give finite real-algebraic upper tests, but no uniform efficient bit-complexity claim is made.
\end{theorem}
\begin{proof}
The prefix of a genuine proper policy satisfies the one-call reserve, so $a_T\le N_{\min}$. Restriction to shorter prefixes proves $a_T$ is nondecreasing. Every prefix completed by the fallback is proper and has exactly the cost $C_T+W_T(H)$; hence $u_T\ge N_{\min}$. Prefixes of the fallback itself give $u_T\le \kappa_0$. Extending an optimizing upper prefix by one fallback call reproduces its exact completed cost, by the identity $H_\theta=g_\theta+\mathcal T_\theta^*H_\theta$ and first-productive commitment. Therefore $u_{T+1}\le u_T$. Compactness of the finite plan space and continuity of its cost operators prove both attainments.

For a lower optimizer, the reserve proof gives
$S_{\theta,T}\preccurlyeq\min\{1,(a_T/c_*-1)/T\}I$. Positivity and $H_\theta\preccurlyeq \kappa_0I$ give $W_{\theta,T}(H_\theta)\preccurlyeq \kappa_0S_{\theta,T}$. Complete that very prefix. Its cost is at most
\[
 a_T+(\kappa_0-c_*)\min\{1,(a_T/c_*-1)/T\}.
\]
Since $a_T\le N_{\min}\le \kappa_0$, this proves \eqref{eq:cost-enclosure}. The squeeze proves both limits. For boundary attainment, the compact product of all finite nodewise decision simplexes has the finite-intersection property for the closed prefix cost constraints at budget $N_{\min}$: feasibility at the largest of finitely many horizons follows from $a_T\le N_{\min}$. A common infinite policy therefore obeys all reserved constraints. Their survival bound tends to zero, and monotone convergence gives properness and cost at most $N_{\min}$. This is the same compactness mechanism as the inherited synthesis theorem, not a finite-memory assertion.

If $N>N_{\min}$, convergence of $u_T$ implies that the strict inequality $u_{T_0}<N$ holds at some finite level. An upper optimizer is itself the finite prefix and a physically executable fallback, not a model-dependent analytic certificate without a common policy. Its start-state uniformity permits it at every later open leaf of a time-homogeneous experiment. Substitute its cost and its counted live labels in the completion and initial-mixture proof; every inequality remains unchanged, except that its $B$ live roles replace the one fallback row. A fixed dummy decision may be replaced by any point of the eventual task dictionary without changing precommit calls or cost. Exact algebraic tests decide a strictly feasible upper threshold when the coefficients and proposed budget have the stated algebraic representation. General real data still require an exact oracle or further numerical certification.
\end{proof}

For example, two orthogonal nonzero failure projections have one-row fair fallback cost $\kappa_0=2$, but fair-first alternation has uniform cost $3/2$ and stops by call two. The cost upper problem at depth two therefore supplies a fallback below every budget $N>3/2$, including $3/2<N<2$. Thus the preliminary optimization removes a genuine restriction of the crude fallback, rather than only renaming its constant.

\subsection{Robustness and scope of the construction}
If two fallback killed maps on the same space satisfy
$\|\mathcal T^*-\widetilde{\mathcal T}^*\|_{\infty\to\infty}\le\epsilon$
on Hermitian matrices, telescoping powers gives
$\|O_d-\widetilde O_d\|\le d\epsilon$: each positive subunital adjoint is a contraction in this norm. Hence $\delta-d\epsilon>0$ certifies a perturbed proper fallback and cost at most $c^*d/(\delta-d\epsilon)$, when the perturbed calls obey the same cost upper bound. This is a quantitative stability certificate for the fallback, not a claim that a terminal-law distance controls unbounded cost. The sharp ladder concerns depth; the scalar rare-success example shows why a spectral margin must enter any numerical certificate.

This section treats finite known candidate models, not statistical learning of an unrestricted unknown kernel. It constructs a common policy without a lower witness, but does not prove that its risk is always the geometric profile $\Q$. The geometric transfer and this primitive synthesis are complementary routes: one gives interpretable sharp resource laws under its certificates, the other supplies a convergent constrained optimum and finite implementation on its finite interface. Quantum reachability, sequence form, potential arguments and semidefinite separation are established ingredients; the present assertion is their explicit paid, same-budget, fixed-output synthesis with a derived fallback and all-model/all-initial-state quantifiers.

\section{Boundary attainment and finite optimality witnesses}\label{sec:primal-dual}
An upper construction and a lower inequality do not by themselves characterize which strategies are optimal. We now close the constrained risk frontier at its feasible boundary, derive its exact variational dual at a fixed dictionary, and identify finite witnesses for each compact-model prefix problem. Cost and task witnesses are allowed to be different physical preparations in different models. Their common-program Bellman certificate is necessary and sufficient at the finite level, and, together with the preceding paid completion, certifies arbitrarily accurate strategies at every interior budget.

The convex-analytic ingredients are classical: finite-support semi-infinite duality is developed in \cite{Charnes63,Shapiro09}, and compactness of strategic and occupation measures has an extensive control-theoretic history \cite{Balder89,Dufour25}. We prove the needed statements in this experiment's coordinates. We do not claim a new general separation or Carath\'eodory theorem. The distinction between a fixed dictionary and a hard bound on the number of freely chosen decoder values is essential: the latter can destroy convexity even in physical dimension one.

\subsection{The exact frontier, including its left endpoint}
For this subsection only, the model family need not be compact, time homogeneous, or finite-dimensional. It suffices that all models share a countable, finitely branching action--report tree with nonempty finite legal action sets at open nodes. Costs have finite conditional means at each finite history, are nonnegative, and have the same positive conditional mean floor $c_*$. Terminal losses lie in $[0,L]$. Initial preparations may be absorbed into the model index. One common observer must commit exactly once and satisfy its expected-cost bound in every model. The decision set $D$ is compact metric and the terminal losses have a common decision modulus. These conditions are satisfied by the primitive interfaces above; they specify the wider scope used here without assuming a free posterior or a common physical reset.

Use the common sequence weights of Lemma~\ref{lem:common-plan}. A fixed ordered decoder tuple $C$ may have repeated entries. Let $\mathcal P_N(C)$ denote the common infinite plans of uniform expected cost at most $N$. Let $q_{C,T}(N)$ be the infimum of worst truncated terminal loss over common depth-$T$ plans with the one-call reserve $C_{\theta,T}+c_*S_{\theta,T}\le N$ in every model and preparation. For an operator interface this inequality means the corresponding Loewner inequality. Define $q_{m,T}$ by optimizing the decoder tuple in $D^m$, and likewise define $R_m$. A finite infeasible problem has value $+\infty$.

\begin{theorem}[Closed constrained frontier and boundary convergence]\label{thm:boundary-attainment}
Every nonempty $\mathcal P_N(C)$ is compact in common-plan coordinates. Terminal-risk functions converge uniformly to their depth-$T$ truncations on this set, with error at most
\[
 L\min\{1,(N/c_*-1)/T\}\qquad(T\ge1).
\]
Thus the optimal risk is attained at every feasible finite budget, including the least feasible budget when one exists. The same assertion holds with a compact decoder tuple of prescribed width $m$.

A budget is feasible if and only if all its finite reserved problems are feasible. At every feasible $N$,
\begin{equation}\label{eq:boundary-convergence}
 q_{C,T}(N)\uparrow R_C(N),\qquad q_{m,T}(N)\uparrow R_m(N).
\end{equation}
Both risk functions are nonincreasing and right-continuous on their feasible budget ray. For a fixed dictionary, $R_C$ is convex and hence continuous at every interior feasible budget. No convexity of $R_m$ is asserted. Boundary convergence gives neither a computable rate nor a finite-memory optimum at the least budget.
\end{theorem}
\begin{proof}
The infinite common-plan space is a closed subset of the countable product of the intervals $[0,1]$, specified by the common flow equalities. It is compact and metrizable. Every prefix cost, survival mass and truncated terminal loss is continuous in its finitely many sequence coordinates for each model. Intersecting all prefix reserve constraints over every model and every depth gives a closed set.

On that intersection, the cost floor and monotonicity of survival give
$C_T+c_*S_T\ge c_*[1+T S_T]$. Hence the displayed uniform survival bound holds, and survival tends to zero. Nonnegative prefix costs increase to the actual cost; their bounds imply a proper common plan of cost at most $N$. Conversely a proper cost-$N$ plan satisfies every reserve because each open path owes a further paid call. This identifies the closed intersection with $\mathcal P_N(C)$, without assuming continuity of total cost or a uniform expected remaining-cost tail.

The nonnegative omitted terminal loss is at most $L$ times survival. A finite prefix loss is uniformly continuous in its finite plan coordinates, uniformly over models: its coefficients are bounded by $L$, and there are finitely many terminal paths. For a decoder tuple, the common decision modulus supplies the additional uniform continuity. Thus the full worst-model risk is a uniform limit of continuous functions on the compact cost-bounded set. It attains its minimum. The product with $D^m$ is compact as well; output indices, not unknown model-dependent values, are the terminal choices.

A later feasible prefix restricts to an earlier one while preserving its reserve and decreasing its accumulated nonnegative loss. Thus the finite values are nondecreasing and bounded above by the actual value. Conversely choose optimizers at depths tending to infinity, extend them by arbitrary legal rows, and take a coordinatewise convergent subsequence, also in $D^m$ when needed. At each fixed earlier depth the limiting reserve and terminal-loss constraints hold in every model. The limit is proper by the preceding argument, and its terminal risk is at most the increasing limit of the finite values. This proves \eqref{eq:boundary-convergence}. The same compact finite-intersection argument without the loss constraints proves the feasibility equivalence. Choosing a sequence of feasible budgets decreasing to their infimum also proves attainment of the least finite feasible cost.

For right-continuity, use optimizers at budgets $N_j\downarrow N$, all bounded by $N+1$, and take a convergent subsequence in that compact cost-bounded set. Its limit is feasible at $N$, and uniform terminal-tail control identifies the limiting risk. Monotonicity supplies the reverse inequality. Finally, two common plans with the same dictionary may be mixed in sequence coordinates. Their costs and losses mix linearly before taking worst models, so the worst cost and risk mix sublinearly. This proves convexity of $R_C$. Such mixing is a common behavioral policy by Lemma~\ref{lem:common-plan}; a finite implementation must count its live branch roles as in the supplement. Mixing two optimized width-$m$ dictionaries can use $2m$ values, so this argument does not prove convexity of $R_m$.
\end{proof}

For a fixed dictionary, write $\mathfrak C(P)$ and $\mathfrak V(P)$ for a plan's worst expected cost and worst terminal risk, and let $\mathcal P_{\rm fin}$ contain the proper common plans of finite uniform cost. The following is a scalar resource dual, not an interchange of a physical model and an observer.
\begin{corollary}[Exact resource-price dual]\label{cor:resource-price}
If some finite budget is feasible, then for every $N\ge N_{\min}$,
\begin{equation}\label{eq:resource-price}
 R_C(N)=\sup_{\lambda\ge0}\inf_{P\in\mathcal P_{\rm fin}}
       \{\mathfrak V(P)+\lambda[\mathfrak C(P)-N]\}.
\end{equation}
At $N>N_{\min}$ the supremum is attained. If a common plan of cost at most $\kappa<N$ is known, an attaining multiplier satisfies $0\le\lambda\le L/(N-\kappa)$. Every primal optimum then minimizes the penalized expression and obeys
$\lambda[\mathfrak C(P)-N]=0$. At $N=N_{\min}$ the supremum remains exact but need not be attained by a finite multiplier. These statements concern one fixed dictionary, not its nonconvex width-constrained union.
\end{corollary}
\begin{proof}
Weak duality follows by testing any feasible plan. On the interior the finite convex nonincreasing function $R_C$ has a supporting line of slope $-\lambda\le0$ at $N$. For every finite-cost plan of cost $n$, its risk is at least $R_C(n)$, which is at least $R_C(N)-\lambda(n-N)$. A primal optimum gives the reverse bound and complementary slackness. Applying the supporting inequality at $\kappa$ gives $\lambda(N-\kappa)\le R_C(\kappa)-R_C(N)\le L$. Supporting lines exist because the left and right secant slopes of a finite convex function are ordered and locally bounded. At the left endpoint choose $N_j\downarrow N_{\min}$ from the interior. The attaining interior multiplier gives a dual value at $N_{\min}$ at least $R_C(N_j)$. Right-continuity and weak duality prove endpoint equality without a bounded-multiplier assumption.
\end{proof}


\subsection{A complete decision comparison of performed stopped laws}
Fix the actual decoder dictionary $C$, rather than only its cardinality. Let $\Omega$ be the measurable disjoint union of terminal action--report words, committed values in $C$, and a prospective audit outcome in a common standard Borel space. The audit kernel after each such word and commitment is fixed in each model. Write $\mathsf P_\theta^P$ for the complete performed stopped law of a legal plan $P$. It includes the stopping call and the raw visible word, not just the final audit marginal. For a specified target family of probability laws $\mathsf T_\theta$ on this same space define
\[
 \delta_{N,C}(\mathsf T)
 =\inf_{P\in\mathcal P_N(C)}\sup_\theta
                \TV(\mathsf P_\theta^P,\mathsf T_\theta),
 \qquad \TV(\mu,\nu)=\sup_A|\mu(A)-\nu(A)|.
\]
The target is data, not an extra executable oracle. Its feasibility in the source interface is the question being tested.

\begin{theorem}[Paid complete-score randomization criterion]\label{thm:complete-stopped-order}
Under the common finite-branching and positive-cost hypotheses of Theorem~\ref{thm:boundary-attainment}, at any feasible finite $N$, the set of families $\{\mathsf P_\theta^P\}_\theta$ is compact and convex in the uniform total-variation metric. Hence the infimum defining $\delta_{N,C}$ is attained, and
\begin{equation}\label{eq:complete-stopped-dual}
 \delta_{N,C}(\mathsf T)
 =\sup_{\pi,g}\left[
 \min_{P\in\mathcal P_N(C)}\sum_{\theta\in\operatorname{supp}\pi}
           \pi_\theta\E_{\mathsf P_\theta^P}g_\theta
 -\sum_{\theta\in\operatorname{supp}\pi}
           \pi_\theta\E_{\mathsf T_\theta}g_\theta\right].
\end{equation}
Here $\pi$ ranges over finitely supported model probabilities and each $g_\theta:\Omega\to[0,1]$ is a finite-valued measurable score. The inner policies remain budget-feasible in \emph{every} model, not only the support of $\pi$. In particular exact common realization of the target at the prescribed budget and dictionary is equivalent to its dominating the optimal source score in every such finite decision test.

On the compact primitive interface, let one common finite-label fallback have conditional cost at most $\kappa<N$. For every $T\ge1$ and every $P\in\mathcal P_N(C)$ there is a common finite-label $P_T\in\mathcal P_N(C)$ satisfying
\begin{equation}\label{eq:stopped-law-density}
 \sup_\theta\TV(\mathsf P_\theta^{P_T},\mathsf P_\theta^P)
 \le D_T+\frac{e_T}{N-\kappa+e_T},\qquad
 D_T=\min\{1,(N/c_*-1)/T\},\quad e_T=(\kappa-c_*)D_T.
\end{equation}
The prior finite-prefix and fallback live-label bounds apply. Therefore finite-label laws are dense at the same budget at every strictly interior feasible budget, and finite-score lower witnesses and actual finite-label realizations approximate the exact deficiency. This is a comparison of the declared performed laws, not of every quantum channel with arbitrary ancillary experiments.
\end{theorem}
\begin{proof}
At depth at most $T$ there are finitely many terminal action--report words and output indices. Their contributions to a stopped law are fixed physical audit measures multiplied by common sequence weights. For two plans, the total-variation difference of these finite contributions is at most the sum of the absolute differences of those weights, uniformly over models, since each physical measure has mass at most one. The two omitted positive tails have mass at most $D_T$ each on the budget set. First choose $T$ large and then make the finitely many coordinates close. This proves continuity of the law-family map into uniform total variation. Compactness of $\mathcal P_N(C)$ gives compactness of its image; linearity in sequence weights gives convexity. Distance to any fixed target family is a continuous function in this metric, so its minimum is attained. This argument does not need continuity of total expected cost or uniform integrability of its tail.

For equal-mass probability measures, $\TV(\mu,\nu)=\sup_A[\mu(A)-\nu(A)]$, since complements reverse the sign. Thus each test functional
$F_{\theta,A}(P)=\mathsf P_\theta^P(A)-\mathsf T_\theta(A)$ is continuous and affine on the compact convex plan set. The right side of \eqref{eq:complete-stopped-dual} is at most the deficiency by the elementary minimax inequality. If $r<\delta_{N,C}$, the closed sets $\{P:F_{\theta,A}(P)\le r\}$ have empty total intersection. Compactness gives a finite subcollection with empty intersection. Its attainable test vectors form a compact convex set disjoint from the lower orthant $(-\infty,r]^q$. Strict separation gives nonnegative weights summing to one, with the weighted test value greater than $r$ for every feasible plan. Nonnegativity follows because a negative coefficient would make the supremum on that lower orthant infinite. Grouping tests with the same model yields a finitely supported $\pi$ and a convex combination of finitely many indicators as $g_\theta$. The latter is finite-valued. This proves the reverse inequality by letting $r$ increase to the deficiency. Zero deficiency implies exact realization by attainment, not only approximation.

For \eqref{eq:stopped-law-density}, keep $P$ unchanged through the depth-$T$ commitment decision, and execute the real fallback on its open paths. Coupling the original and completed experiments with the same calls and current coins until then makes their complete stopped laws identical unless the original path is still open. The total-variation discrepancy is therefore at most $D_T$, including all raw history and audit coordinates. The completed cost is at most $N+e_T$, because its reserved prefix already paid $c_*$ for each open path. Mix this completion with the common fallback with probability $p=(N-\kappa)/(N-\kappa+e_T)$. The expected cost is at most $N$ in every model. Mixing changes a probability law by at most $1-p$ in total variation. The triangle inequality proves the displayed bound. Its finite implementation counts the actual finite tree, the fallback and the initial branch role; the same dictionary is used throughout. The compact-family cost theorem provides such a fallback at every $N>N_{\min}$. Applying this approximation to an attaining deficiency minimizer proves the final assertion.
\end{proof}

The decision-separation principle belongs to the Blackwell--Le Cam tradition \cite{Blackwell,Jencova16}. The additional work here is the compactness of the \emph{performed} law family under a uniform paid stopping constraint, together with finite-label approximation without increasing that constraint. It requires all bounded measurable scores of that performed outcome; localized affine-square tests alone do not suffice. The conclusion does not grant a free stateful simulator, an uncharged audit on rejected calls, or control of unbounded expected duration by total variation. The latter is controlled by the separate budget throughout. If the dictionary itself varies, laws recording different numerical readouts need not vary continuously in total variation; the general-loss width theorem and the fixed-dictionary comparison are distinct statements.

\subsection{Sparse spectral certificates on compact uncertainty}
Return to the continuous compact finite-dimensional interface of Section~\ref{sec:compact-risk}. Fix a dictionary and horizon. Let $\mathcal X_T$ be the common finite sequence-plan polytope and let $k_T$ be its affine dimension, after eliminating flow equalities. The same statement applies to a common tree with compulsory terminal nodes, or to a prescribed face of its plan polytope; $k_T$ then denotes the affine dimension of that polytope. For $z=(\theta,u)$ with $\theta\in\Theta$ and $\|u\|=1$, put
\[
 c_z(x)=\langle u,(C_{\theta,T}(x)+c_*S_{\theta,T}(x))u\rangle,
 \qquad v_z(x)=\langle u,V_{\theta,T}(x)u\rangle.
\]
These are continuous families of affine functions of $x$; $0\le v_z(x)\le L$. Testing all such $z$ is equivalent to the operator constraints. Write $q_T(N)$ for the resulting compact-model finite optimum.

\begin{theorem}[Exact finite witness and complementary slackness]\label{thm:sparse-witness}
Suppose some $x_0\in\mathcal X_T$ satisfies $c_z(x_0)\le\kappa<N$ for every $z$. Then there are risk atoms $z_j$, cost atoms $w_i$, and nonnegative weights $\beta_j,\alpha_i$ such that
\begin{align}
 &\sum_j\beta_j=1,\qquad
 \lambda:=\sum_i\alpha_i\le\frac{L}{N-\kappa},\qquad
 \#\{j:\beta_j>0\}+\#\{i:\alpha_i>0\}\le k_T+1,\label{eq:sparse-size}\\
 &q_T(N)=\min_{x\in\mathcal X_T}
 \left\{\sum_j\beta_jv_{z_j}(x)+
                    \sum_i\alpha_i[c_{w_i}(x)-N]\right\}.\label{eq:sparse-dual}
\end{align}
For any fixed optimal plan $x_*$ these atoms can be chosen active there:
\[
 \beta_j>0\Rightarrow v_{z_j}(x_*)=q_T(N),\qquad
 \alpha_i>0\Rightarrow c_{w_i}(x_*)=N.
\]
Their preparations are pure; active vectors are respectively eigenvectors of the loss operator at $q_T$ and the reserved-cost operator at $N$. Conversely, a feasible plan, such active atoms and minimization of the displayed affine expression certify optimality. The total atom bound is sharp for each affine dimension $k_T$. It concerns dual witnesses, not the number of controller labels or a fixed horizon-independent adversarial prior.
\end{theorem}
\begin{proof}
Set $v(x)=\max_zv_z(x)$ and $c(x)=\max_zc_z(x)$. Compactness gives attained maxima. The finite value as a function of its budget is convex and nonincreasing. A supporting line at $N$, which is interior because $\kappa<N$, yields a multiplier $\lambda\ge0$ for which $x_*$ minimizes $v+\lambda(c-N)$ and $\lambda(c(x_*)-N)=0$, exactly as in Corollary~\ref{cor:resource-price}.

For completeness, the finite polytope argument underlying its stationarity is proved here. The directional derivative of a compact maximum of affine functions is the maximum of their active gradients: take a maximizing index along a sequence of positive steps, pass to a convergent subsequence, and compare with every active index. Let $A$ be the compact convex sum of the active risk-gradient hull and $\lambda$ times the active cost-gradient hull; omit the second term if $\lambda=0$. If $A$ did not intersect the negative normal cone of $\mathcal X_T$ at $x_*$, separation of a compact convex set and a closed convex cone would give a tangent direction $h$ with $\max_{a\in A}a\cdot h<0$. For a polytope, every vector in this tangent cone gives a feasible sufficiently short step: impose its finitely many active linear inequalities, then choose the step short enough for the inactive ones. The directional derivative would be negative, a contradiction. Thus there are active finite convex combinations whose sum $g$ satisfies $g\cdot(x-x_*)\ge0$ for all $x\in\mathcal X_T$. Convex combinations in a finite-dimensional compact hull can be chosen finite by eliminating affine dependencies.

In the $k_T$-dimensional affine coordinates, append one coordinate to the active gradients. Represent $(g,1)$ as a conic sum of $(\nabla v_z,1)$ for risk atoms and $(\nabla c_w,0)$ for cost atoms. If more than $k_T+1$ generators have positive coefficients, a linear dependence permits subtracting a suitable multiple of that dependence until one coefficient vanishes and none becomes negative. Repeat. This preserves $g$ and risk mass one, and all remaining atoms stay active. Their affine combination is minimized at $x_*$. Its value there is $q_T$, proving \eqref{eq:sparse-dual}. Evaluating that minimum at $x_0$ gives $q_T\le L-\lambda(N-\kappa)$, hence the bound in \eqref{eq:sparse-size}, even if compression changed the original multiplier.

For sufficiency, at any feasible $x$ the expression in braces is at most $v(x)$, while at the specified active minimizing plan it equals its risk. This sandwiches the optimum. If a positive semidefinite matrix $NI-C$ has zero quadratic form at a unit vector, it annihilates that vector; the spectral claims follow, and similarly for risk. Finally take a compulsory one-call scalar experiment with $k+1$ actions and $k+1$ models, unit cost, one decoder and loss $1$ exactly when action and model indices coincide. With commitment required after that call, its plan polytope is the $k$-simplex and its risk is $\max_i x_i$. At any $N>1$, cost atoms cannot be active. The value $1/(k+1)$ requires the risk weights all equal $1/(k+1)$, since the dual minimum is their smallest weight. Every one of the $k+1$ atoms is necessary. This proves sharpness within physical experiments, not just an abstract cone.
\end{proof}

Finite-support duality and the conic reduction are classical \cite[Theorem 4.4]{Shapiro09}. The operational content is their use with the actual common plan, the open-path reserve, distinct spectral cost and score preparations, and a completion at the same resource budget. An arbitrary coarse grid is not a replacement for this certificate: active models may lie in the interior of a parameter set, as Example~\ref{ex:vertex-risk} shows.

\begin{proposition}[Common Bellman evaluation and equality of used rows]\label{prop:dual-bellman}
For any finite atom lists with risk mass one, the right side of \eqref{eq:sparse-dual} is evaluated exactly by a backward recursion on the single action--report tree. At each productive report it minimizes between a decision in the fixed dictionary and continuation; nonproductive reports must continue. The open depth-$T$ cost is the actual one-call reserve. There is a deterministic finite plan minimizing this dual expression, but it need not satisfy the original robust budget and risk constraints. A randomized plan minimizes it if and only if every choice of positive sequence weight has zero backward slack at its reached node.
\end{proposition}
\begin{proof}
Propagate, for each atom, its unnormalized physical state along each raw action--report word, without multiplying by controller probabilities. At a pre-call word $h$, the immediate cost for action $a$ is $\sum_i\alpha_i\tr(G_{w_i,a}\rho_{w_i,h})$. At its productive report $y$, stopping with $r$ contributes $\sum_j\beta_j\tr(F_{z_j,a,y,r}\rho_{z_j,hay})$. Continuation contributes the next node's recursively computed value. At an open leaf of depth $T$, assign $c_*\sum_i\alpha_i\tr\rho_{w_i,h}$. Add over reports and minimize over the common action. Subtract $\lambda N$ once at the root. This evaluates the desired affine functional by backward induction; the full formula is recorded in the supplement.

Each minimization admits a deterministic choice because its alternatives are finite. Keeping the chosen current command or report never reveals the model: the propagated matrices are offline analytic coefficients for a common tree. At every row, its randomized value minus its minimum is a nonnegative slack. Backward substitution expresses the root excess as the sum of those slacks times the corresponding nonnegative sequence weights. It vanishes precisely when every positively used alternative has zero slack. A dual-minimizing plan alone may violate the individual model constraints; Theorem~\ref{thm:sparse-witness} requires their separate primal verification.
\end{proof}

\begin{corollary}[Finite certificates for the full stopped optimum]\label{cor:complete-certificate}
Fix $N>N_{\min}$ and a dictionary $C$ on the compact primitive interface. For every $\eta>0$ there is a finite horizon, at most $k_T+1$ spectral witness atoms, and a common proper finite-label observer at cost at most $N$ such that the Bellman-evaluated dual lower and that observer's worst risk differ by at most $\eta$. The same actual dictionary is used at every history and by the fallback. With compact optimized output width, attainment still holds, but this dual is conditional on the dictionary; global certification requires the separate dictionary optimization or a certified decision net. No minimax swap over dictionaries is authorized.
\end{corollary}
\begin{proof}
The compact-family cost hierarchy supplies a finite fallback of uniform cost $\kappa<N$. Its prefix is strictly feasible for every reserved finite problem. Apply Theorem~\ref{thm:sparse-witness}. The compact-model problem imposes all model constraints, so completion and initial mixing introduce no model-cover error. With $D_T=\min\{1,(N/c_*-1)/T\}$ and $e_T=(\kappa-c_*)D_T$, its realized risk is at most
$q_T+LD_T+Le_T/(N-\kappa+e_T)$. This difference tends to zero, while every dual value lower-bounds every feasible full-history strategy. Choose $T$ to make the difference at most $\eta$. The prior completion proof counts the finite prefix, fallback and initial-mixture roles. For variable dictionaries use the preceding decision-net or algebraic optimization, not convexification of their union.
\end{proof}


\begin{theorem}[Simultaneous expected resources]\label{thm:multi-resource}
The compact-model certificate and completion extend to finitely many additive nonnegative resource costs. Write $C_{\theta,T}^{(\ell)}$ for the prefix cost of resource $\ell=0,\ldots,r$, assume conditional lower and upper bounds $0\le c_\ell I\preccurlyeq G_{\theta,a}^{(\ell)}\preccurlyeq c_\ell^*I$, and require $c_0>0$. Fix a budget vector $\mathbf N$ and one common finite fallback with uniform resource bounds $\kappa_\ell<N_\ell$ simultaneously. Each finite constraint is
\[
 C_{\theta,T}^{(\ell)}+c_\ell S_{\theta,T}
                         \preccurlyeq N_\ell I.
\]
Every finite optimal plan admits an exact certificate with risk mass one and cost atoms indexed by a resource, a model and a pure preparation. If $\lambda_\ell$ is the total multiplier mass for resource $\ell$, then
\[
 \sum_{\ell=0}^r\lambda_\ell(N_\ell-\kappa_\ell)\le L,
 \qquad \text{total positive cost and risk atoms}\le k_T+1.
\]
Thus the support bound depends on the affine plan dimension, not directly on the number of resource constraints. The common Bellman recursion adds all weighted immediate resource costs; its reserve sums $c_\ell$ times the weighted open mass of the cost atoms for resource $\ell$.

Let $q_T(\mathbf N)$ denote the finite lower and put
\[
 D_T=\min\{1,(N_0/c_0-1)/T\},\quad
 e_{\ell,T}=(\kappa_\ell-c_\ell)D_T,\quad
 p_T=\min_{0\le\ell\le r}
 \frac{N_\ell-\kappa_\ell}{N_\ell-\kappa_\ell+e_{\ell,T}}.
\]
There is a common proper finite-label observer satisfying every original budget and having worst loss at most
\[
 q_T(\mathbf N)+LD_T+L(1-p_T).
\]
The difference tends to zero. At any feasible finite vector, primal attainment and monotone finite lower convergence also hold, without asserting a finite boundary implementation. The costs may represent separately charged calls, time, energy or communication, when they obey these additive hypotheses; a hard program-size, scratch-space or dictionary-cardinality constraint is not such a cost by declaration.
\end{theorem}
\begin{proof}
For the finite problem set $c_\ell(x)=\max_{\theta,u}\langle u,(C_{\theta,T}^{(\ell)}+c_\ell S_{\theta,T})u\rangle$ and $v(x)=\max_{\theta,u}\langle u,V_{\theta,T}u\rangle$. The upward closed set of achievable vectors $(c_0(x)-N_0,\ldots,c_r(x)-N_r,v(x))$ with coordinatewise slack is convex and closed, since $x$ lies in a compact polytope. A supporting hyperplane at $(0,\ldots,0,q_T)$ has nonnegative coefficients. Its risk coefficient is positive: if it were zero, the jointly strictly feasible $x_0$ would give a strictly smaller pairing. Normalize that coefficient to one. This yields multipliers $\lambda_\ell\ge0$, minimization of $v+\sum_\ell\lambda_\ell(c_\ell-N_\ell)$ and complementary slackness.

Repeat the active-gradient and normal-cone argument with one cost-gradient hull for each positive multiplier. All cost generators still have appended coordinate zero, and risk generators have appended coordinate one. The conic reduction in dimension $k_T+1$ therefore proves the same total support bound. Evaluating the resulting affine minimum at the jointly strict fallback prefix gives the displayed weighted margin bound. At each node the immediate scalar cost is simply the sum over resource-indexed atoms, which proves the recursion claim.

Resource 0 bounds survival by $D_T$. Completing an optimizer by the one common fallback costs at most $N_\ell+e_{\ell,T}$ in resource $\ell$ and adds at most $LD_T$ risk. Choose this completed plan with probability $p_T$ and the fallback otherwise. For every $\ell$ its cost is at most
$p_T(N_\ell+e_{\ell,T})+(1-p_T)\kappa_\ell\le N_\ell$.
The risk bound follows exactly as for one resource. Strict positive margins make $p_T\to1$. Common-tree and fallback labels are counted as before. On a general feasible vector, intersect all prefix reserve constraints; the single coercive resource 0 forces properness and uniform bounded-loss tails. The compactness and finite-intersection proof of Theorem~\ref{thm:boundary-attainment} then applies with all the additional closed constraints. Separate fallbacks safe for separate resources would not justify this proof: joint feasibility of one common fallback is essential.
\end{proof}

\subsection{Distinct witnesses and the price of hard output width}
\begin{example}[A shared adversarial prior can give a strict dual gap]\label{ex:distinct-dual-witnesses}
There are two commands $A,B$, one productive report, one decoder, and two scalar models. Costs are charged but are not extra observations; the visible action--report history has the same law in both models. Model 1 has costs $(1,2)$ and losses $(0,0)$; model 2 has costs $(1,1)$ and losses $(1,0)$. Commitment may occur only after a call. At budget $N=3/2$, writing $x$ for the probability that the final command is $B$, every plan costs at least $1+x$ in model 1 and has risk $1-x$ in model 2. One-call mixing at $x=1/2$ attains the exact value $1/2$. The dual certificate has one unit cost atom at model 1 and one unit risk atom at model 2.

If one incorrectly forces both atom lists to be proportional to the same probability $(t,1-t)$, the one-call dual becomes
\[
 \sup_{0\le t\le1,\lambda\ge0}\min_{0\le x\le1}
 \{(1-t)(1-x)+\lambda(tx-1/2)\}
       =\tfrac32-\sqrt2<\tfrac12.
\]
For $t\le1/2$ its best value is zero; for $t>1/2$ the best multiplier is $(1-t)/t$, giving $3/2-t-1/(2t)$, maximized at $t=1/\sqrt2$. All quantities come from real cost and loss tables. A single least-favorable prior for one scalarized task is not a substitute for independently imposed robust cost and risk constraints.
\end{example}

\begin{proposition}[An exact hard-width convexification gap]\label{prop:width-dual-gap}
There is a one-model, one-dimensional physical experiment with two commands, compact decision set $D=\{0,1\}$ and bounded nonnegative loss, such that for $1\le N\le2$,
\[
 R_1(N)=\min\{1/2,2-N\},\qquad R_2(N)=(2-N)/2.
\]
Thus $R_1$ is attained and right-continuous but not convex. Its scalar Lagrange relaxation over all width-one programs equals $R_2$; at $N=3/2$ it has a gap $1/4$. Randomizing two different one-value dictionaries pays for a two-value dictionary, not for an uncharged initial seed.
\end{proposition}
\begin{proof}
Both commands always give one productive report, with costs $1,2$. Decision 0 has loss $1/2$ after either command. Decision 1 has loss $1$ after the cheap command and zero after the costly one. These are permissible prospective audit losses: the future audit can be the deterministic command symbol with the displayed loss table. This is a bounded general-loss example, not a claim about a squared-loss realization.

For the sole decoder 0 every proper program has loss $1/2$. For the sole decoder 1, if its final command is costly with probability $x$, expected total cost is at least $1+x$ and risk is $1-x$. One-call mixing realizes all these pairs, so optimization gives $R_1$. With both values, the smallest loss is $1/2$ after a cheap last command and zero after a costly one. The same cost lower and one-call realization give $R_2$.

Over width-one programs the penalized infimum is
$\min\{1/2+\lambda(1-N),\lambda(2-N)\}$: the cheap decoder-1 endpoint is dominated. Maximizing over $\lambda\ge0$ gives $(2-N)/2$, with equality of the two lines at $\lambda=1/2$. At $N=3/2$ the primal width-one value is $1/2$ and the relaxation is $1/4$. In the randomized implementation attaining the latter, a cheap call uses decoder 0 and a costly call uses decoder 1, so two distinct final values are genuinely required.
\end{proof}

The preceding results supply an exact variational and finite-witness completion of the primitive risk route. They do not establish that every causal family admits the geometric reference certificates, a common finite unknown-model predictive quotient, or optimal computational resources. Boundary primal attainment is compatible with the failure of uniform expected cost-tail convergence and with lack of a finite boundary implementation. Those independent questions are not resolved by a success certificate for a finite optimization problem.

\section{Weighted boundary decisions and positive stochastic realization}\label{sec:robust-realization}
The deterministic residual quotient of Section~\ref{sec:boundary-memory} does not classify a robust problem with tied Bayes actions. In that case the object to realize is a probability-valued decision function, and its residuals need not take finitely many values. We first derive that object without assuming model-independent report probabilities. We then characterize its finite stochastic implementation by an invariant convex family, and give exact memory classifications for a random-horizon calibration task. The convex realization machinery is classical in origin: linear representations and positive realizations are different notions \cite{Jaeger00,ThonJaeger15,MonrasWinter14}. The new assertions concern the physical minimum-cost face, its robust optimal decision face, and their counted stochastic implementation.

\subsection{State-dependent reports and a minimum-cost potential}
Fix a compact metric model space $\Theta$, a finite-dimensional physical space $K$, a nonempty finite failure alphabet $Y$, and a fixed nonempty finite decoder dictionary $C$. All primitive matrices below depend continuously on $\theta$. The prepared density is $\rho_\theta$. There is one globally legal command, with completely positive branches $\mathcal S_\theta$ for a productive herald and $\mathcal K_{\theta,y}$ for failure $y$, whose sum is trace preserving. A rejected herald retains the actual post-$\mathcal S_\theta$ state; a failure retains the actual post-$\mathcal K_{\theta,y}$ state. No replacement preparation, inspection of the protected state, or additional unpriced operation is available. Calls have nonnegative cost with mean observable
\[
 c_* I\preccurlyeq G_\theta\preccurlyeq c^*I,\qquad 0<c_*\le c^*<\infty.
\]
Commitment is allowed only on a herald and precedes the prospective audit. Its loss at decoder $r$ has observable $0\preccurlyeq F_{\theta,r}\preccurlyeq LI$ on the post-herald state. An independently executed final audit cost must be added separately. The declared reports are the whole observer input; charged costs do not supply an additional observation channel. These are single-command hypotheses, not a finite-memory criterion for a general controlled POMDP.

Put $\mathcal T_\theta=\sum_y\mathcal K_{\theta,y}$ and $d=\dim K$. Assume
\begin{equation}\label{eq:weighted-depth}
 I-(\mathcal T_\theta^*)^dI\succ0\qquad(\theta\in\Theta).
\end{equation}
This is the primitive depth test already proved in Theorem~\ref{thm:primitive-fallback}, now applied to the one command. Compactness gives a common $\delta>0$ with the left side at least $\delta I$. Thus first-herald survival after $n$ calls is at most
\[
 b_n=(1-\delta)^{\lfloor n/d\rfloor}.
\]
The uniformly convergent resolvent
\[
 H_\theta=(I-\mathcal T_\theta^*)^{-1}G_\theta,
 \qquad c_*I\preccurlyeq H_\theta\preccurlyeq K_0I,
 \quad K_0=c^*d/\delta
\]
is continuous in $\theta$. Write $h_\theta=\tr(H_\theta\rho_\theta)$, $N_0=\max_\theta h_\theta$ and $\Theta_0=\{\theta:h_\theta=N_0\}$. For $w=y_1\cdots y_n$, let $\mathcal K_{\theta,w}=\mathcal K_{\theta,y_n}\cdots\mathcal K_{\theta,y_1}$ and
\[
 \lambda_\theta(w)=\tr\mathcal S_\theta\mathcal K_{\theta,w}(\rho_\theta).
\]
The empty word is included. Words may have zero probability. We use the following support condition:
\begin{equation}\label{eq:cost-support}
 \lambda_\theta(w)>0\text{ for some }\theta
 \quad\Longrightarrow\quad
 \lambda_{\theta'}(w)>0\text{ for some }\theta'\in\Theta_0.
\end{equation}
This is a sufficient coverage condition for robust boundary rigidity, not a claimed necessary condition for every such rigidity. In particular it holds if one cost-maximizing model has $\mathcal S^*I\succ0$ and $\mathcal K_y^*I\succ0$ for every $y$. These inequalities ensure that every finite word has positive mass in that model; neither its posterior nor the model label is supplied to the observer.

\begin{theorem}[Weighted herald rigidity]\label{thm:weighted-rigidity}
For every proper finite-cost observer in a fixed model,
\begin{equation}\label{eq:weighted-residual}
 \mathbb E_\theta C_\tau
 =h_\theta+
 \mathbb E_\theta\sum_{i\le\tau:\,\text{herald }i\text{ rejected}}
            \tr(H_\theta\rho_i^+),
\end{equation}
where $\rho_i^+$ is its normalized post-herald physical state. Consequently the least robust budget is $N_0$. Under \eqref{eq:cost-support}, every budget-$N_0$ observer commits at the first herald in every model. Model- and state-dependent failure and herald probabilities are permitted.
\end{theorem}
\begin{proof}
At a reached pre-call state $\rho$, the expected cost plus the potential after continuing failures is
$\tr[(G_\theta+\mathcal T_\theta^*H_\theta)\rho]=\tr(H_\theta\rho)$.
Continuing after a herald adds exactly its post-herald potential, weighted by the actual rejection decision. Conditional summation to a deterministic horizon gives the displayed identity with cost and rejected-herald sums truncated, minus the remaining live potential. That remainder is bounded by $K_0\Pr_\theta(\tau>n)$ and tends to zero. Finite cost and the positive cost floor imply proper survival decay; monotone convergence applies to the two nonnegative sums. This proves \eqref{eq:weighted-residual}. First-herald commitment realizes $h_\theta$ simultaneously, so its largest value is the least robust budget.

At every $\theta'\in\Theta_0$, a boundary observer has zero residual. Since $H_{\theta'}\succeq c_*I$, it rejects no herald with positive probability. Before the first herald, the distribution of its private memory conditional on the failure word depends only on that word and its own coins: there is one command, and the physical report maps do not depend on those coins. In particular its rejection probability at the first herald after $w$ is model independent. If $w$ is possible anywhere, \eqref{eq:cost-support} makes it possible at a cost-maximizing model, where this probability is zero. Thus the probability is zero in every model. The sufficient full-support test follows by successively taking traces of positive, nonzero branch outputs.
\end{proof}

The coverage qualification cannot be dropped. With unit calls, let one model have geometric herald rate $1/3$ and only failure letter $a$. Let another have a prepared two-state system whose first call fails with letter $b$, enters a state with deterministic heralds, and thereafter remains there. Both pass \eqref{eq:weighted-depth} (use a two-dimensional identity extension in the first model). Their first-herald costs are three and two. A policy accepting all $a$-histories but rejecting the first herald after $b$ still has robust cost three. The word $b$ is absent from the expensive model. This failure of coverage licenses an actual extra stopping decision; it is not a zero-probability technicality.

\subsection{The optimal face with arbitrary robust ties}
Assume \eqref{eq:cost-support}. Define the unnormalized loss coefficients
\[
 f_\theta(w,r)=\tr\bigl(F_{\theta,r}\mathcal S_\theta
                  \mathcal K_{\theta,w}(\rho_\theta)\bigr),
 \quad 0\le f_\theta(w,r)\le L\lambda_\theta(w).
\]
A boundary rule is $b=(b_w)_{w\in Y^*}\in\mathcal D:=\prod_{w\in Y^*}\Delta(C)$, independent of $\theta$, with risk
$R_\theta(b)=\sum_{w,r}b_w(r)f_\theta(w,r)$.
For a probability measure $\mu$ on $\Theta$ put
$f_\mu(w,r)=\int f_\theta(w,r)\,d\mu(\theta)$ and
$A_\mu(w)=\arg\min_{r\in C} f_\mu(w,r)$.

\begin{theorem}[Weighted minimax face and tied optimality]\label{thm:weighted-face}
The full-history boundary value and a least-favorable measure are attained, and
\begin{equation}\label{eq:weighted-minimax}
 V=\min_{b\in\mathcal D}\max_\theta R_\theta(b)
   =\max_{\mu\in\mathcal P(\Theta)}
           \sum_w\min_{r\in C}f_\mu(w,r).
\end{equation}
For any attaining $\mu_*$ the complete optimal face is
\begin{equation}\label{eq:optimal-face}
 \mathcal B_*=
 \{b\in\mathcal D:R_\theta(b)\le V\ (\theta\in\Theta),\quad
               \supp b_w\subseteq A_{\mu_*}(w)\ (w\in Y^*)\}.
\end{equation}
Words with zero $\mu_*$-averaged mass impose no extra decoder restriction. In general the support restrictions alone are insufficient: the inequalities for the individual models remain essential.
\end{theorem}
\begin{proof}
The terms beyond length $n-1$ have sum at most $Lb_n$, uniformly in $b$ and $\theta$. Finite partial sums are jointly continuous, and this bound proves joint continuity on the compact convex rule space and on the compact model space. Passing from models to their probability measures makes the payoff bilinear and continuous. Compact minimax gives \eqref{eq:weighted-minimax}, with both extrema attained; equivalently apply finite-dimensional separating minimax to finite word truncations and then use the uniform tail and compactness. Minimization for fixed $\mu$ is wordwise, since each coordinate ranges over a simplex.

For an optimal $b$, its $\mu_*$-average is both at most $V$ and at least the Bayes value $V$. Subtraction yields the nonnegative identity
\[
 0=\sum_{w,r}b_w(r)
       [f_{\mu_*}(w,r)-\min_u f_{\mu_*}(w,u)].
\]
Every summand vanishes. This gives the support restriction even without positive model-independent word masses; if a whole word has zero averaged mass, every coefficient there is zero and all actions are allowed. Conversely the model inequalities make $b$ robustly optimal, since $V$ is the infimum. This proves the full face, not sufficiency of Bayes ties alone. For example two constant loss vectors $(0,1)$ and $(1,0)$ have fair robust optimum $1/2$, while both pure actions are tied under the equal model mixture and each has worst risk one.
\end{proof}

\subsection{What replaces a deterministic residual congruence}
View each $g\in\mathcal D$ as a probability-valued response function on future words. Its left shift is $(L_yg)(z)=g(yz)$. An \emph{anchored finite convex realization} of $b$ consists of $g_1=b,g_2,\ldots,g_M\in\mathcal D$ with
\begin{equation}\label{eq:convex-invariance}
 L_yg_s\in\operatorname{conv}\{g_1,\ldots,g_M\}
       \quad(s=1,\ldots,M,\ y\in Y).
\end{equation}
The distinguished initial response $g_1$ is counted even if it is not an extreme point of this convex hull. This anchor reflects a prescribed initial live label; it must not be replaced by a free random initial state. The auxiliary $g_s$ need not themselves be robust optimal from a new preparation.

\begin{theorem}[Exact stochastic boundary memory with ties]\label{thm:convex-memory}
Under the weighted rigidity hypotheses, an exact robust boundary optimum with at most $M$ live labels exists if and only if some $b\in\mathcal B_*$ admits an anchored convex realization with at most $M$ generators. The least such generator count is the least exact live-label count. This criterion permits all robust ties, state-dependent word masses and stochastic transitions.

For a specified response $b$, the block Hankel matrix
\[
 \mathsf H_b(w,(z,r))=b_{wz}(r)
\]
has rank at most the generator count. Finite Hankel rank alone is not sufficient. Event-stopped implementations use at most $\max(M,|C|)$ total labels, and state-stopped implementations at most $M+|C|$; unused decoder values can be removed from these upper counts.
\end{theorem}
\begin{proof}
From \eqref{eq:convex-invariance} choose stochastic coefficients $Q_y(s,t)$ representing each shifted generator. Set the decoder row $D_s=g_s(\varnothing)$. Induction on word length gives
\[
 g_s(y_1\cdots y_n)=e_s^{\mathsf T}Q_{y_1}\cdots Q_{y_n}D.
\]
The machine starts in label one, updates on failures with these matrices, and on the first herald draws its decoder from $D_s$. It uses no physical-state input and realizes $b$ exactly, so \eqref{eq:optimal-face} proves optimality.

Conversely, for each reached label of any stochastic machine define $g_s(z)$ to be its terminal decoder distribution after processing $z$ and then a herald. Its rows are probabilities. The Markov property gives $L_yg_s=\sum_tQ_y(s,t)g_t$, and the initial function is $b$. Unreachable labels may be removed; any unused rows can be completed arbitrarily. Thus the same number of reached labels gives the required generators, proving both directions and minimality. Factor the Hankel matrix as the row of reached state probabilities after $w$ times the column of future response probabilities from each state after $z$ to obtain the rank lower bound. The established event/state role packing gives the total-label upper counts.
\end{proof}

This is a positive realization criterion on the optimal decision face, not a claim to originate residual automata or finite-dimensional linear realization. Rabin--Scott and the Moore/Myhill--Nerode tradition characterize deterministic finite residual functions \cite{RabinScott59}; observable-operator and multiplicity-automaton theories separate linear rank from stochastic realization \cite{Jaeger00,ThonJaeger15}. Here exact unique Bayes actions recover the deterministic criterion of Section~\ref{sec:boundary-memory}. With ties, a convex family can have infinitely many different residual points, and selecting positive successors no longer preserves the required output probabilities. The next application exhibits both phenomena exactly.

For finite primitive models the criterion also has a direct operator verification form. Given stochastic rows $Q_y,D$, define the killed map on the direct sum of $M$ physical blocks by
\[
 (\mathcal T_{\theta,Q}X)_t
      =\sum_{s,y}Q_y(s,t)\mathcal K_{\theta,y}(X_s),
 \qquad A_{\theta,s}=\sum_rD_s(r)\mathcal S_\theta^*F_{\theta,r}.
\]
The physical survival bound also bounds this map, independently of the memory rows. Its unique score potential is
\begin{equation}\label{eq:positive-score-certificate}
 Z_\theta-\mathcal T_{\theta,Q}^*Z_\theta=A_\theta,
 \qquad 0\preccurlyeq Z_\theta\preccurlyeq LI.
\end{equation}
Its initial risk is $\tr((Z_\theta)_1\rho_\theta)$. Consequently, when an exact target value $V$ is supplied, stochastic rows satisfying this risk at most $V$ in every model are a finite necessary-and-sufficient certificate for exact $M$-label optimality. The equations follow by summing the absolutely convergent killed score series; normalization bounds its sum by $LI$, and transience proves uniqueness. With finitely many algebraic models and algebraic $V$, this is a finite real-algebraic existence test. It does not assume that the unknown optimum of an arbitrary problem is itself algebraic.

\subsection{Robust Born-probability replication at an unknown random horizon}
Consider a protected system prepared in $\rho$, one trace-preserving channel $\Phi$, and an effect $0\preccurlyeq E\preccurlyeq I$. On each call a failure has probability $t$ and applies $\Phi$; a herald has probability $1-t$ and leaves the system unchanged. The fixed unknown parameter ranges over $[t_-,t_+]\subset(0,1)$, with $t_-<t_+$. Calls cost one. On commitment output $r\in C=\{0,1\}$ before observing the Bernoulli audit $Z$ with effect $E$. There are two audit-loss models $\sigma\in\{-1,1\}$, neither affecting report probabilities, with loss
\begin{equation}\label{eq:calibration-loss}
 \ell_\sigma(r,Z)=\tfrac12[1+\sigma(r-Z)]\in[0,1].
\end{equation}
The unknown sign enforces two-sided calibration rather than choosing one favorable signed discrepancy. This is a replication task for the Born probability at a random horizon, not the squared-sign prediction task of the previous qubit theorem. The observer has no access to the protected system except the failure/herald reports.

Write $\beta_n=\tr(E\Phi^n\rho)$ and let $b_n$ be the probability of output one after $n$ failures. The endpoint $t_+$ has the largest cost and full word support, so Theorem~\ref{thm:weighted-rigidity} forces first-herald commitment at the least budget $N_0=1/(1-t_+)$. A boundary policy has
\begin{equation}\label{eq:analytic-calibration}
 \sup_{t,\sigma}R_{t,\sigma}(b)
 =\tfrac12+\tfrac12\sup_{t\in[t_-,t_+]}
       \left|(1-t)\sum_{n\ge0}t^n(b_n-\beta_n)\right|.
\end{equation}

\begin{theorem}[Unique randomized optimum and exact realization memory]\label{thm:calibration-realization}
The full-history boundary value is $1/2$, and its unique response kernel is $b_n=\beta_n$ for all $n\ge0$. An exact finite-memory optimum exists precisely when
\begin{equation}\label{eq:born-realization}
 \beta_n=e_1^{\mathsf T}Q^n d\quad(n\ge0)
\end{equation}
for a finite row-stochastic matrix $Q$ and a vector $d\in[0,1]^M$; its least size is the exact live memory. Under a least-favorable measure putting equal weight on the two signs at any fixed $t$, every wordwise action is Bayes-tied. Thus uniqueness of the robust randomized kernel does not imply uniqueness of any wordwise Bayes action.

Put $g(t)=(1-t)\sum_{n\ge0}t^n\beta_n$. The exact one-live-label value is
\begin{equation}\label{eq:calibration-one}
 V_1=\tfrac12+\tfrac14(\max_tg(t)-\min_tg(t)).
\end{equation}
For every fixed $M$ the best risk is attained. Whenever \eqref{eq:born-realization} is impossible at that size, its risk is strictly greater than $1/2$. A $T$-label unary prefix chain gives a boundary-feasible risk at most $1/2+t_+^T/2$.
\end{theorem}
\begin{proof}
A direct expectation of \eqref{eq:calibration-loss} gives \eqref{eq:analytic-calibration}. The choice $b_n=\beta_n$ proves value $1/2$, which is also a lower bound from the two opposite signs. At equality the power series $\sum_n(b_n-\beta_n)t^n$ vanishes on a nonempty interval inside its unit disk of convergence. The analytic identity theorem makes all its coefficients zero. For completeness, a nonzero convergent power series is locally its first nonzero power times a nonvanishing analytic factor at any interior zero; hence its zeros cannot accumulate in its disk unless it is identically zero. This proves uniqueness.

Theorem~\ref{thm:convex-memory} in the unary alphabet gives exactly \eqref{eq:born-realization}. Averaging the two loss signs makes every conditional loss $1/2$, proving the least-favorable tie assertion. A one-label rule is a constant output probability $u$. Its discrepancy is $u-g(t)$; the midpoint of the minimum and maximum of $g$ minimizes the uniform absolute error, with radius half their difference. This proves \eqref{eq:calibration-one}.

For fixed $M$, the compact set of stochastic rows and decoder values has continuous risk: each finite sum is continuous and the tail is uniformly at most $t_+^T$. Attainment and unique optimality prove the strict gap. Store labels $0,\ldots,T-1$, advance until the final label and repeat it thereafter, with outputs $\beta_0,\ldots,\beta_{T-1}$. It is exact for $n<T$ and differs by at most one later. Formula~\eqref{eq:analytic-calibration} bounds its risk excess by $t_+^T/2$. Its stopping is unchanged, so this is the same least budget, not an enlarged-budget approximation.
\end{proof}

\begin{theorem}[Every finite exact memory dimension occurs]\label{thm:calibration-k}
For every $k\ge2$ and algebraic $a\in(0,1)$ there is a $k$-state protected classical channel in this task whose least exact stochastic live memory is exactly $k$. Take an initial chain of $k-1$ transient states: each stays with probability $a$ and advances with probability $1-a$; state $k$ is absorbing. Audit whether the state is still transient. Then
\begin{equation}\label{eq:phase-survival}
 \beta_n=\sum_{j=0}^{k-2}\binom nj(1-a)^j a^{n-j},\qquad
 g(t)=1-\left(\frac{(1-a)t}{1-at}\right)^{k-1}.
\end{equation}
The exact one-label excess is $[g(t_-)-g(t_+)]/4$. For $k=2$, infinitely many different probabilistic residuals are implemented by just two stochastic labels.
\end{theorem}
\begin{proof}
The number of advances before time $n$ is binomial until absorption, proving the formula for $\beta_n$ with $\binom nj=0$ for $j>n$. An independent observer copy of the transition probabilities, using only private coins and $k$ labels, gives the required decoder probability; it neither reads nor copies the actual protected state. This proves the upper memory bound.

As a function of $n$, $\beta_n=a^nP(n)$ for a polynomial $P$ of degree exactly $k-2$. The shifted polynomials $P(n),\ldots,P(n+k-2)$ span all polynomials of degree at most $k-2$: successive finite differences lower the degree by exactly one, with nonzero leading coefficient. The scalar Hankel row/column rank is therefore $k-1$. Including both output probabilities $\beta_n,1-\beta_n$ adds the constant function, which is independent because every exponential-polynomial term tends to zero. Thus the probability-valued Hankel rank is $k$, giving the matching lower by Theorem~\ref{thm:convex-memory}. Summing the binomial series, using $\sum_{n\ge j}\binom nj x^{n-j}=(1-x)^{-j-1}$, gives $g$. It decreases with $t$, giving the stated one-label formula. For $k=2$ the residuals have first coordinate $a^n$, all distinct, although the displayed chain has two labels.
\end{proof}

\subsection{A sharp positive-realization transition for the protected qubit}
Use the algebraic rotation from the preceding section,
\[
 U=5^{-1/2}\begin{pmatrix}2&-1\\1&2\end{pmatrix},\quad
 c=\cos\vartheta=3/5,\ s=\sin\vartheta=4/5,
\]
with $\rho=|0\rangle\langle0|$, $E=|0\rangle\langle0|$ and
$\Phi_\alpha(A)=\alpha UAU^*+(1-\alpha)\tr(A)I/2$, $0\le\alpha\le1$.
For each fixed known retention factor $\alpha$, the required response is
\begin{equation}\label{eq:damped-born}
 \beta_n=\tfrac12[1+\alpha^n\cos(n\vartheta)],
\end{equation}
with $\beta_0=1$ also at $\alpha=0$.

\begin{theorem}[Exact qubit realization phases]\label{thm:qubit-positive-phase}
For the robust calibration task, the following hold.

\textup{(a)} At $\alpha=1$ no finite classical memory attains the boundary optimum, although the probability-valued Hankel rank is three. At $\alpha=0$ the least exact live memory is two.

\textup{(b)} At every $0<\alpha<1$ an exact stochastic optimum exists with at most
\[
 M_\alpha=\max\{3,\lceil\pi/\arccos\alpha\rceil\}
\]
live labels. The least number tends to infinity as $\alpha\uparrow1$.

\textup{(c)} Put $\alpha_3=5/(3+4\sqrt3)$. The least exact live memory is three precisely when $0<\alpha\le\alpha_3$. It is at least four for $\alpha_3<\alpha<1$. The three-label assertion is global over all stochastic transition matrices and randomized decoder rows, not only circulant candidates.
\end{theorem}
\begin{proof}
For $\alpha>0$, the three distinct modes $1$, $\alpha e^{i\vartheta}$ and $\alpha e^{-i\vartheta}$ have nonzero coefficients in \eqref{eq:damped-born}. Vandermonde independence gives Hankel rank three. Cayley--Hamilton applied to any representation \eqref{eq:born-realization} shows that its transition matrix must have the two latter modes as eigenvalues: substituting its characteristic recurrence and using independence of the three exponential sequences makes the characteristic polynomial vanish at each mode.

Every unit-modulus eigenvalue $z$ of a finite stochastic matrix is a root of unity. Indeed choose an eigenvector coordinate of maximal modulus. Equality in the triangle inequality for its stochastic row forces every positive-probability successor to have the same maximal modulus and its value to equal $z$ times the current value. Follow positive successors to a repeated state. The directed cycle gives $z^j=1$ for some positive integer $j$. But $e^{i\vartheta}$ is not a root of unity: otherwise $2c=6/5$ would be a rational algebraic integer and hence an integer. This proves (a) at $\alpha=1$. At $\alpha=0$ a starting label with output one followed by an absorbing fair-output label realizes the sequence, whereas one label cannot change its initial output.

For (b), take the regular $M_\alpha$-gon with vertices $v_j$ on the unit circle, including $v_0=(1,0)$. Its inradius is at least $\alpha$, so $\alpha R_\vartheta$ maps every vertex into the polygon. Choose barycentric coefficients $Q(j,l)$ of each image in these vertices and let $d_j=(1+(v_j)_1)/2$. Then $Q$ is stochastic, all $d_j\in[0,1]$, and the expected vertex after $n$ transitions from vertex zero is $(\alpha R_\vartheta)^nv_0$. This proves \eqref{eq:damped-born} exactly. If least exact memories stayed bounded along $\alpha\uparrow1$, compactness of their padded transition/decoder rows would give, by taking a subsequence and each finite coefficient limit, a finite realization at $\alpha=1$. This contradicts (a).

For (c), define a circulant stochastic matrix with first row
\[
 u_0=\frac{1+2\alpha c}{3},\qquad
 u_+=\frac{1-\alpha c+\sqrt3\alpha s}{3},\qquad
 u_-=\frac{1-\alpha c-\sqrt3\alpha s}{3}.
\]
Its other rows are the cyclic translates. All entries are nonnegative exactly up to the asserted threshold. Its complex nonconstant eigenvalues are $\alpha(c\pm is)$, and decoder $(1,1/4,1/4)$ from the first state realizes \eqref{eq:damped-born}. The rank lower makes three minimal.

Conversely any three-state realization has eigenvalues $1,\alpha(c+is),\alpha(c-is)$. For $B=I-Q$ put $a_i=B_{ii}=\sum_{j\ne i}Q_{ij}\ge0$. Its two nonzero eigenvalues are $u\pm iv$, where $u=1-\alpha c$ and $v=\alpha s$. Its trace and second characteristic coefficient give
\[
 \sum_i a_i=2u,\qquad
 u^2+v^2=\sum_{i<j}(a_ia_j-Q_{ij}Q_{ji})
       \le\sum_{i<j}a_ia_j\le\tfrac13(2u)^2.
\]
Therefore $\sqrt3\alpha s\le1-\alpha c$, exactly $\alpha\le\alpha_3$. No assumption of irreducibility, symmetry or deterministic decoding entered this lower bound.
\end{proof}

The damped channel has no strict limiting preference between the two signed calibration losses, and its exact optimal decoder law remains nonconstant at arbitrarily late times. Positive stochastic realization nevertheless succeeds for every $\alpha<1$. This does not contradict the earlier impossibility for unique squared-sign prediction: the fixed losses, the robust optimal face and the output law differ. The unit-circle obstruction and invariant positive cones belong to classical positive-realization theory \cite{MonrasWinter14}; the task-specific exact count, the global three-label threshold and the same-boundary physical implementation are proved here. No sharp asymptotic memory rate is inferred from the polygon upper.

\begin{proposition}[Finite algebraic certificates and minimal synthesis]\label{prop:finite-realization-test}
Suppose a probability-valued response has a supplied real-algebraic linear representation
$b_w=a^{\mathsf T}T_wB$ of dimension $r$ on a finite alphabet. Existence of an exact stochastic realization with $M$ live labels is equivalent to the finite conditions
\[
 Q_y\ge0,\ Q_y\mathbf1=\mathbf1,\quad D\ge0,\ D\mathbf1=\mathbf1,
 \qquad e_1^{\mathsf T}Q_wD=a^{\mathsf T}T_wB
       \quad(|w|\le M+r-1).
\]
These are real-algebraic feasibility conditions, hence decidable at each prescribed $M$, and a feasible instance has algebraic controller entries. For the qubit family with algebraic $0<\alpha<1$, the finite upper from Theorem~\ref{thm:qubit-positive-phase} makes successive size tests a terminating synthesis of a minimal exact stochastic observer. No efficient bit bound or general decision procedure for unbounded positive realization is asserted.
\end{proposition}
\begin{proof}
Use the block diagonal representation with transition $\operatorname{diag}(Q_y,T_y)$, initial row $(e_1^{\mathsf T},-a^{\mathsf T})$ and final columns $(D,B)^{\mathsf T}$. Its reachable row spaces grow from dimension at most one in ambient dimension $M+r$. Equality of successive spaces implies invariance under every letter, so the space generated by words of length at most $M+r-1$ contains all word rows. Annihilating the final columns on these finitely many words therefore annihilates them on every word. The reverse implication is immediate. The conditions are polynomial in the controller entries, with real-algebraic coefficients, so real-closed-field elimination decides feasibility and yields an algebraic sample when nonempty \cite{Seidenberg54}. The polygon construction supplies a finite upper; one can obtain a certified integer bound by testing $\alpha<\cos(\pi/M)$ with increasing integers, using algebraic comparisons. Thus finitely many prescribed-size tests find the first feasible size. This is a complete terminating construction for this family, not a claim about the complexity of general probabilistic automata.
\end{proof}

\section{Finite-memory optimality at the least feasible budget}\label{sec:boundary-memory}
Boundary attainment in Theorem~\ref{thm:boundary-attainment} concerns the compact space of full-history plans. It does not assert that the optimum belongs to any finite-label class. We resolve that distinction on a class of prepared heralded channel experiments. The result is a necessary-and-sufficient memory criterion, an explicit finite-dimensional quantum obstruction, and quantitative approximation at the same boundary budget. It is not an additional assumption of the general compactness theorem.

\subsection{Prepared heralded channels and a boundary rigidity identity}
Let $0<p<1$, put $q=1-p$, and fix a nonempty finite failure alphabet $Y$ with probabilities $\pi_y>0$ summing to one. There is one globally legal unit-cost command. Model $\theta$ has a prepared density matrix $\rho_\theta$ on a fixed finite-dimensional space and channels $\Phi_{\theta,y}$. Its report instrument is
\begin{equation}\label{eq:herald-instrument}
 \mathcal I_{\theta,s}(A)=pA,\qquad
 \mathcal I_{\theta,y}(A)=q\pi_y\Phi_{\theta,y}(A)\quad(y\in Y).
\end{equation}
The herald $s$ is productive; a failure $y$ forbids commitment. Rejection of $s$ leaves the physical state unchanged. The observer cannot inspect the protected physical state. After commitment to $r$ in one fixed nonempty finite dictionary $C$, the future audit has conditional loss observable $0\preccurlyeq F_{\theta,r}\preccurlyeq LI$. Any separately charged fixed final-audit cost is added to all budgets below. No unpriced channel, phase input or clock is supplied to the observer.

For $w=y_1\cdots y_n\in Y^*$, including the empty word, set
\[
 \rho_\theta(w)=\Phi_{\theta,y_n}\cdots\Phi_{\theta,y_1}(\rho_\theta),\quad
 \lambda(w)=pq^{|w|}\prod_{i=1}^{|w|}\pi_{y_i},\quad
 \ell_\theta(w,r)=\tr(F_{\theta,r}\rho_\theta(w)).
\]
Every word has positive mass, $\sum_w\lambda(w)=1$, and $\sum_{|w|\ge T}\lambda(w)=q^T$. Probabilities of visible failures are independent of both the hidden model and the private controller coins. The channels themselves may change the prepared state noncommutatively.

\begin{lemma}[The least budget forces first-herald commitment]\label{lem:boundary-herald}
The least feasible precommit expected-call budget is $N_{\min}=1/p$. For every proper finite-cost policy, let $J$ count productive reports which it rejects. Then, in every model,
\begin{equation}\label{eq:herald-cost-identity}
 p\E_\theta\tau=1+\E_\theta J.
\end{equation}
Consequently a policy is feasible at $N_{\min}$ if and only if it commits at the first herald, up to null histories. At that budget the only choice is its decision distribution $b_w\in\Delta(C)$ after each failure word $w$, and its risk is
\begin{equation}\label{eq:boundary-risk-functional}
 \mathcal R_\theta(b)=\sum_w\lambda(w)\sum_{r\in C}b_w(r)\ell_\theta(w,r).
\end{equation}
\end{lemma}
\begin{proof}
The event that call $i$ is made is predictable, and the conditional herald probability at that call is $p$. Summing the nonnegative indicators gives $\E_\theta\sum_{i\le\tau}\one_{\{s_i=s\}}=p\E_\theta\tau$. Exactly one of these heralds is accepted; all others are rejected. This proves the identity without an optional-stopping integrability assumption. It implies the lower $1/p$, attained by first-herald commitment. Equality forces the nonnegative integer $J$ to vanish almost surely. Conversely that rule has a geometric call count of mean $1/p$. Conditional on a failure word, the physical state is the displayed channel product and is independent of the private tape. The independent, unobserved audit is performed only after the decision. This gives the risk formula.
\end{proof}

\subsection{A robust decision quotient and its exact memory content}
Assume the model space $\Theta$ is compact metric and the prepared states, channels and audit observables are continuous in $\theta$. A singleton is permitted; an arbitrary prior is not part of the data. Write $V=\min_b\sup_\theta\mathcal R_\theta(b)$ at $N_{\min}$. For a probability measure $\mu$ on $\Theta$ put
\[
 \ell_\mu(w,r)=\int\ell_\theta(w,r)\,d\mu(\theta),\quad
 a_\mu(w)=\min_{r\in C}\ell_\mu(w,r),\quad
 V(\mu)=\sum_w\lambda(w)a_\mu(w).
\]

\begin{theorem}[Boundary minimax and the regularity criterion]\label{thm:boundary-regularity}
The primal optimum and a least-favorable measure are attained, and
\begin{equation}\label{eq:boundary-minimax}
 V=\max_{\mu\in\mathcal P(\Theta)}V(\mu).
\end{equation}
For an attaining $\mu_*$, every robust optimum has zero expected decision slack:
\begin{equation}\label{eq:boundary-slack}
 \int\mathcal R_\theta(b)\,d\mu_*(\theta)-V
 =\sum_w\lambda(w)\sum_r b_w(r)
          [\ell_{\mu_*}(w,r)-a_{\mu_*}(w)]\ge0.
\end{equation}
Suppose the minimizer $a(w)\in C$ of $\ell_{\mu_*}(w,\cdot)$ is unique for every word. Then $a$ is the unique robust optimal decision rule, and a finite stochastic autonomous observer attains $V$ if and only if the relation
\begin{equation}\label{eq:decision-congruence}
 w\equiv_a v\quad\Longleftrightarrow\quad
          a(wz)=a(vz)\text{ for every }z\in Y^*
\end{equation}
has finite index. The least number of live labels attaining $V$ is exactly that index. Random transitions and randomized readouts cannot reduce this exact count.

At each prescribed finite live-label bound $M$, the best boundary risk is attained. If the index is infinite, that best risk is strictly larger than $V$ for every finite $M$, although finite-label policies converge to $V$ at the same budget. A prefix implementation uses at most
$1+\sum_{j=0}^{T-1}|Y|^j$ live labels and has excess risk at most $Lq^T$. Its output dictionary remains $C$. For any implementation with $M$ live labels and at most $m$ actual decoder values, event stopping uses at most $\max(M,m)$ total labels and state stopping uses at most $M+m$, by Lemma~\ref{lem:roles}.
\end{theorem}
\begin{proof}
The product $\mathcal B=\prod_{w\in Y^*}\Delta(C)$ is compact and convex. The tail of every risk series is at most $Lq^T$, uniformly in $b$ and $\theta$. Each truncated expression is jointly continuous and affine in the appropriate variable. The full payoff on $\mathcal B\times\mathcal P(\Theta)$ is therefore continuous and bilinear, with both factors compact convex. The elementary compact minimax theorem applies \cite{Sion58}. For completeness, it can also be obtained here from the finite-prefix games: truncate at $T$, separate the convex hull of the finite-dimensional model loss vectors from a strict lower orthant, and take compact subsequential limits of priors and plans. The uniform error $Lq^T$ passes both inequalities to the limit. Minimizing for a fixed measure chooses a minimizing decision separately at every word and gives $V(\mu)$. This proves equality and both attainments. Nonnegative summation proves the slack identity. At a robust optimum its left side is at most zero and hence is zero. Since every $\lambda(w)>0$, uniqueness forces $b_w$ to be the point mass at $a(w)$. Existence of a robust optimum then proves that this Bayes rule is itself robustly optimal; a general Bayes minimizer need not have this property when ties occur.

At the boundary an $M$-label stochastic observer is represented by one stochastic transition matrix $Q_y$ for each failure letter and one decoder distribution $d_s\in\Delta(C)$ at each live label. Its initial live label is fixed. Unreachable rows can be completed arbitrarily. Zero slack says that every label reached with positive probability after $w$ decodes $a(w)$ with probability one. Choose one positive-probability successor at each pair $(s,y)$ and the uniquely prescribed output on each reached label. This gives a deterministic observer with no more labels and the same output at every word: its path always had positive probability in the stochastic observer. Thus exact random realization implies deterministic realization.

If a deterministic machine sends $w,v$ to the same label, every common continuation sends them through identical transitions and gives the same output. Hence they are equivalent in \eqref{eq:decision-congruence}. Any exact machine has at least its number of equivalence classes. Conversely this relation is a right congruence: $w\equiv_a v$ implies $wy\equiv_a vy$. Its classes, with transition $[w]\mapsto[wy]$ and decoder $a(w)$, form an exact deterministic observer whenever the index is finite. This is the standard residual-function construction for finite automata, applied to the optimal decision rule rather than to the full predictive state.

For fixed $M$, all $Q_y$ and $d_s$ lie in a finite product of compact simplexes. Each finite risk sum is continuous in their entries. The same uniform geometric tail makes full risks, and their supremum over $\Theta$, continuous. A minimum is attained. It can equal $V$ only if an exact machine exists, proving the strict gap for infinite index. Finally follow an optimal full-history rule on words of length less than $T$, storing that finite tree, and use one fallback decoder label thereafter. The policy still stops at the first herald and costs exactly $1/p$. Only the terminal mass $q^T$ can change loss, proving the approximation and state bound. This argument uses no interior budget slack.
\end{proof}

The quotient in \eqref{eq:decision-congruence} is specific to this boundary task and this fixed dictionary. It is not a common predictive quotient preserving every experiment or future audit. The minimax argument is classical; the new content is its combination with exact minimum-cost commitment, zero-slack determinization, and a necessary-and-sufficient autonomous memory count. Finite-index residual characterizations and unary periodicity are established automata theory \cite{Mereghetti24}; their use does not establish a new priority claim for those tools.

\begin{proposition}[Ties for one preparation, and why robust ties differ]\label{prop:boundary-ties}
For a single prepared model, without uniqueness, an exact finite stochastic observer exists if and only if there is a finite-index right congruence $\sim$ on $Y^*$ such that
\[
 \bigcap_{w\in A}\arg\min_{r\in C}\ell(w,r)\ne\varnothing
             \quad\text{for every equivalence class }A.
\]
The least number of live labels is the least number of classes among such congruences. For several unknown models the same zero-slack condition is necessary for an optimal stochastic observer under a least-favorable measure, but deterministic selection from tied Bayes actions need not preserve robust optimality.
\end{proposition}
\begin{proof}
For one preparation, every wordwise minimizing rule is globally optimal. Zero slack makes the decoder support of any reached stochastic label a subset of every optimal-action set at its reaching words. Choose one such output and a positive successor for each failure letter. Determinization preserves optimality. The partition of words by its reached deterministic label is a right congruence with the displayed nonempty intersections. Conversely choose an output in each intersection and run the congruence-class observer. Minimizing the number of classes proves the count.

For the distinction, take two models with identical channels and reports and losses $(0,1)$ and $(1,0)$ on a two-point dictionary, independent of the word. The robust value is $1/2$, attained with one live label and a fair terminal decision. Under the equally weighted least-favorable measure both actions are tied at loss $1/2$. Either constant deterministic selection has worst loss one. Thus Bayes zero slack alone cannot license robust determinization in the presence of ties.
\end{proof}

\begin{proposition}[A finite distinguishing-depth lower certificate]\label{prop:boundary-depth-gap}
In the unique-action case of Theorem~\ref{thm:boundary-regularity}, suppose the index is infinite and $m=|C|\ge2$. Fix $M$. There is a finite $H_M$ such that every deterministic $M$-label observer disagrees with $a$ at a word of length at most $H_M$. Put
\[
 \gamma_H=\min_{|w|\le H,\ r\ne a(w)}
       [\ell_{\mu_*}(w,r)-\ell_{\mu_*}(w,a(w))]>0,
 \qquad \pi_{\min}=\min_y\pi_y.
\]
Every stochastic $M$-label boundary observer obeys
\begin{equation}\label{eq:boundary-quantitative-gap}
 \sup_\theta\mathcal R_\theta-V
       \ge\frac{p\gamma_{H_M}}{m}
                    \left(\frac{q\pi_{\min}}{M}\right)^{H_M}>0.
\end{equation}
This is a finite distinguishing certificate, not a sharp asymptotic formula for the optimal intermediate-memory risk.
\end{proposition}
\begin{proof}
There are finitely many deterministic transition and output tables on $M$ named labels. Each table fails at some finite word because the residual index is infinite. The maximum of their first distinguishing depths is finite. From any stochastic table choose a largest-probability successor for each $(s,y)$, of probability at least $1/M$, and a largest-probability output at each label, of probability at least $1/m$. The resulting deterministic table fails at a word $w$ of length $n\le H_M$. In the stochastic table the selected internal path and output occur with probability at least $M^{-n}/m$ conditional on this word. Their slack is at least $\gamma_{H_M}$. The word's terminal mass is at least $p(q\pi_{\min})^n$. Keep just that term in \eqref{eq:boundary-slack}; since $q\pi_{\min}/M<1$, lowering it to the exponent $H_M$ gives the bound. Worst-model risk is no smaller than the least-favorable average.
\end{proof}

\subsection{An algebraic qubit with no finite-memory boundary optimum}
A quantum orbit supplies a finite-dimensional physical realization, not an arbitrary nonregular string placed in an external clock. Take one failure letter and a single prepared qubit,
\begin{equation}\label{eq:qubit-boundary-data}
 \rho_0=|0\rangle\langle0|,\qquad
 U=\frac1{\sqrt5}\begin{pmatrix}2&-1\\1&2\end{pmatrix},\qquad
 \Phi(A)=UAU^*,\qquad C=\{-1,1\}.
\end{equation}
On commitment the future audit measures $Z=\operatorname{diag}(1,-1)$, with outcomes $z\in\{-1,1\}$, and scores $(z-r)^2$. Thus $F_r=2I-2rZ$ and $L=4$. Write $\vartheta$ for the Bloch rotation angle, so $\cos\vartheta=3/5$, $\sin\vartheta=4/5$, and put $c_n=\cos(n\vartheta)$. In particular $c_0=1$, $c_1=3/5$, and $c_{n+1}=(6/5)c_n-c_{n-1}$.

\begin{theorem}[Boundary nonattainment by every finite classical memory]\label{thm:qubit-boundary}
For every $0<p<1$, the least budget is $1/p$ and the full-history risk there is
\begin{equation}\label{eq:qubit-boundary-value}
 V_p=2-2p\sum_{n=0}^\infty q^n|c_n|.
\end{equation}
Its unique optimal rule is $a_n=\operatorname{sgn}c_n$. It has infinite residual index. Hence no finite stochastic autonomous observer attains $V_p$ at this budget. This remains true with arbitrary real transition probabilities and randomized terminal decisions.

Let $V_{p,M}$ be the optimal risk with $M$ live labels and the same two-point dictionary. Each such minimum exists, and
\begin{equation}\label{eq:qubit-memory-gap}
 0<V_{p,M}-V_p\le4q^M.
\end{equation}
The exact one-live-label optimum is
\begin{equation}\label{eq:qubit-one-label}
 V_{p,1}=2-\frac{2p(1-3q/5)}{1-6q/5+q^2}.
\end{equation}
At $p=1/2$ it is $12/13$, while the full-history value is the strictly smaller convergent series in \eqref{eq:qubit-boundary-value}. All specified primitive entries are algebraic. The bound \eqref{eq:boundary-quantitative-gap} is an explicit positive lower with $m=2$, $\pi_{\min}=1$, and $\gamma_H=4\min_{0\le n\le H}|c_n|$.
\end{theorem}
\begin{proof}
After $n$ failures the prepared state is $U^n\rho_0U^{*n}$, whose audit mean is $c_n$. The expected squared loss of decision $r$ is $2-2rc_n$. Lemma~\ref{lem:boundary-herald} forces first-herald commitment, so independent pointwise minimization gives \eqref{eq:qubit-boundary-value}.

The number $e^{i\vartheta}$ is not a root of unity. Otherwise $2\cos\vartheta=6/5$ would be a rational algebraic integer, hence an integer, a contradiction. Thus $\vartheta/\pi$ is irrational; no $c_n$ is zero, and each residue subsequence $(n_0+jr)\vartheta$ is dense modulo $2\pi$ for every positive integer $r$. The latter elementary rotation fact follows, for example, by the pigeonhole approximation of a nonzero arbitrarily small rotation and its successive multiples. Each residue subsequence therefore visits both open semicircles on which cosine has opposite signs. The optimal sign sequence cannot be ultimately periodic. A finite deterministic unary observer always follows a finite initial path and then a cycle, so it produces an ultimately periodic decision sequence. Theorem~\ref{thm:boundary-regularity} now excludes every finite stochastic observer and proves the strict gap and its attainment at fixed $M$.

Use $M$ live labels to store $n=0,\ldots,M-1$, incrementing on failure until the last label, which thereafter repeats. Decode the correct sign at these $M$ labels and commit at the first herald. The rule is exact for $n<M$ and its excess loss is bounded by four on the tail of mass $q^M$. Its cost is still exactly $1/p$. This proves the upper in \eqref{eq:qubit-memory-gap}.

For one live label the expected output is one constant $b\in[-1,1]$. Its risk is $2-2bA$, where
\[
 A=p\sum_{n\ge0}q^nc_n
   =p\operatorname{Re}\frac1{1-qe^{i\vartheta}}
   =\frac{p(1-3q/5)}{1-6q/5+q^2}>0.
\]
Thus $b=1$ is optimal and gives \eqref{eq:qubit-one-label}. Substituting $p=q=1/2$ gives $12/13$. The failure and herald maps are genuine completely positive trace-nonincreasing maps with trace-preserving sum. They use one declared command and a protected terminal audit; the proof never reads that audit early or turns the qubit into observer-owned memory.
\end{proof}

\begin{corollary}[A compact uncertain-preparation family]\label{cor:polarization-boundary}
Replace the known initial qubit by
$\rho_\theta=(I+\theta Z)/2$, with unknown fixed $\theta\in[\eta,1]$ and $0<\eta\le1$. Keep the same instrument, dictionary and audit. The robust boundary value is
\[
 V_{p,[\eta,1]}=2-2\eta p\sum_{n\ge0}q^n|c_n|.
\]
A least-favorable preparation is $\rho_\eta$, and the unique robust optimal rule is again $\operatorname{sgn}c_n$. It is unattainable by every finite stochastic autonomous memory, while the same-budget prefix approximation remains valid. Thus the obstruction is not tied to a chosen Bayesian prior or to one pure initial preparation. It concerns this declared preparation family, not all qubit density matrices.
\end{corollary}
\begin{proof}
At preparation $\theta$ the conditional mean is $\theta c_n$. The sign rule minimizes loss word by word for every positive $\theta$, and its total risk decreases with $\theta$. Its worst value is therefore the displayed one at $\eta$. Every common policy is lower-bounded by its risk in that same preparation, whose unique minimizer is the sign rule. This proves the robust value, the least-favorable measure $\delta_\eta$, and uniqueness without an assumed prior. Apply Theorem~\ref{thm:boundary-regularity}. If $\eta<1$, the family includes both mixed and pure states; no inference is made at $\eta=0$, where the least-favorable audit is a fair bit and uniqueness fails.
\end{proof}

\begin{corollary}[Exact classification on unitary orbits]\label{cor:unitary-period}
For any finite-dimensional unitary failure channel, fixed prepared state, fixed finite dictionary and audit observables, suppose the minimizing decision along the orbit is unique at every nonnegative integer time. A finite-label boundary optimum exists if and only if the optimal decision sequence is periodic from time zero. Its least live-label count is the least positive period.
\end{corollary}
\begin{proof}
The residual theorem first gives equivalence with ultimate periodicity. If the sequence has eventual period $r$, there are integers $k_j\to\infty$ with $U^{rk_j}\to I$ in operator norm, by compactness of the finite-dimensional unitary torus. For each fixed $n$, the states at times $n+rk_j$ and $n+r+rk_j$ converge to those at $n$ and $n+r$. Uniqueness among finitely many decisions gives a strict gap at each of these two limiting states, so their optimal decisions eventually agree with the limiting decisions. Eventual $r$-periodicity makes these decisions equal. Therefore $a_n=a_{n+r}$ for every $n\ge0$. A cycle of length the least period realizes the rule, and every finite unary machine has an eventual cycle at least that long.
\end{proof}

\subsection{Dissipation and the strict decision-gap condition}
Finite memory is not confined to periodic unitary orbits. A stable decision at the limiting state gives a different positive realization mechanism.

\begin{proposition}[Uniform contraction with a strict limiting decision gap]\label{prop:contractive-boundary}
Consider a single prepared heralded model with finitely many failure letters. Suppose there is a density matrix $\rho_\infty$ and constants $C_0<\infty$, $0<\alpha<1$ such that
\[
 \|\rho(w)-\rho_\infty\|_1\le C_0\alpha^{|w|}
                \quad(w\in Y^*).
\]
Suppose one dictionary value $r_*$ has a strict limiting gap
\[
 g=\min_{r\ne r_*}\tr[(F_r-F_{r_*})\rho_\infty]>0.
\]
For any nonnegative integer $T$ with $2LC_0\alpha^T<g$, there is an exact optimal boundary observer with at most
$1+\sum_{n=0}^{T-1}|Y|^n$ live labels, where the empty sum is zero. All words of length at least $T$ use the same decision $r_*$. Ties at earlier words are harmless in this single-model statement. In particular, dissipative channels with a certified contraction and a strict limiting task gap can have an exact finite-memory optimum at the least feasible cost.
\end{proposition}
\begin{proof}
For every $r\ne r_*$ and every word of length $n$, trace duality gives
\[
 \ell(w,r)-\ell(w,r_*)
 \ge g-\|F_r-F_{r_*}\|_\infty
                    \|\rho(w)-\rho_\infty\|_1
 \ge g-2LC_0\alpha^n.
\]
This is strictly positive at $n\ge T$. Store all shorter words in a finite tree and merge their depth-$T$ continuations into one label decoding $r_*$. Choose a least-loss value separately at each shorter word. Every decision minimizes its actual conditional loss, and the rule stops at the first herald, so it attains the single-model optimum at cost $1/p$. The displayed number counts the entire prefix tree and the common tail label. No limiting state is supplied to the observer during execution; it is an offline certificate of its finite program.
\end{proof}

The limiting gap is a real restriction, not a consequence of mixing. The following family keeps a strictly positive contraction rate but retains the infinite exact-memory obstruction.
\begin{corollary}[Noisy rotating qubits: contraction is not sufficient]\label{cor:noisy-boundary}
In Theorem~\ref{thm:qubit-boundary}, replace the failure channel by
\[
 \Phi_\alpha(A)=\alpha UAU^*+(1-\alpha)\tr(A)I/2,
                    \qquad 0\le\alpha\le1.
\]
For $0<\alpha\le1$, the exact boundary value and one-live-label value are
\begin{align*}
 V_{p,\alpha}&=2-2p\sum_{n\ge0}(q\alpha)^n|c_n|,\\
 V_{p,\alpha,1}&=2-\frac{2p(1-3q\alpha/5)}{1-6q\alpha/5+q^2\alpha^2}.
\end{align*}
No finite stochastic observer attains $V_{p,\alpha}$, while
\[
 0<V_{p,\alpha,M}-V_{p,\alpha}
       \le \frac{4p(q\alpha)^M}{1-q\alpha}.
\]
For $0<\alpha<1$ the failure channel is uniformly trace-norm contractive to $I/2$. At $\alpha=0$, however, one live label is exactly optimal, with value $2-2p$. Thus a positive mixing margin alone cannot replace the strict limiting decision gap in Proposition~\ref{prop:contractive-boundary}.
\end{corollary}
\begin{proof}
Depolarization commutes with unitary conjugation, so
$\rho_n=I/2+\alpha^n U^n(\rho_0-I/2)U^{*n}$ and its audit mean is $\alpha^nc_n$. The unique minimizing decision is therefore the same nonperiodic sign sequence whenever $\alpha>0$. Apply Theorem~\ref{thm:boundary-regularity}. A wrong decision at time $n$ has excess loss at most $4\alpha^n$; summing the tail from $M$ gives the displayed bound. For a constant decision, replace $q$ by $q\alpha$ in the elementary cosine series, retaining the factor $p$, to obtain the one-label value. Trace differences of states have zero trace and contract by exactly the factor $\alpha$ under this channel. At $\alpha=0$ every positive-length failure word has state $I/2$ and both decisions are tied, while at the empty word the decision $+1$ is optimal. The constant $+1$ rule attains $2-2p$. At $I/2$ the two limiting losses agree, explaining why the preceding positive-gap theorem does not apply.
\end{proof}

\subsection{The exact two-label frontier for a noisy quantum orbit}
The qubit application also admits a complete nontrivial finite-memory optimization. The next formula allows every rotation angle and damping strength, rather than fixing the algebraic example's angle.

\begin{theorem}[Exact two-live-label quantum prediction law]\label{thm:two-label-quantum}
Consider the same prepared qubit, herald and squared Pauli audit, with a Bloch rotation angle $\vartheta$, $c=\cos\vartheta\in[-1,1]$, and depolarizing strength $0\le\alpha\le1$. Put $x=q\alpha$ and
\[
 f_c(t)=\frac{1-ct}{1-2ct+t^2},\qquad A=p f_c(x).
\]
The optimal one-live-label risk is $2-2A$. The exact optimum among all stochastic observers with at most two live labels is $2-2S_2$, where the following cases determine $S_2$ and an attaining transition matrix. On a herald the first label decodes $+1$ and the second $-1$, except in the constant-decision cases.

If $x=0$, then $S_2=p$ and one label suffices. For $x>0$ and $0\le c\le1$, let $s=\sqrt{1-c^2}$ and $t_*=c/(1+s)$. Then
\[
 S_2=\begin{cases}
 A,&x\le t_*,\\
 p(1+s)/s-A,&x>t_*.
 \end{cases}
\]
The second case never occurs at $c=1$. In the first case use a constant $+1$ decision. In the second case an optimal failure matrix is
\[
 Q=\begin{pmatrix}t_*/x&1-t_*/x\\0&1\end{pmatrix}.
\]
For $x>0$ and $-1\le c<0$, write $k=-c$, $s=\sqrt{1-k^2}$, and
\[
 t=\min\left\{x,\max\left\{0,
       \frac{x+k-s\sqrt{1+2kx+x^2}}{1+kx}\right\}\right\}.
\]
Then
\[
 S_2=A+\frac{2[p f_c(-t)-A]}{1+t/x},\qquad
 Q=\begin{pmatrix}0&1\\t/x&1-t/x\end{pmatrix}
\]
is attaining. All formulas include ties and degenerate transitions. The first compulsory call is included in the cost, which remains exactly $1/p$ in every case.

In particular, for the irrational algebraic orbit with $c=3/5$, $p=1/2$ and $\alpha=1$, the exact two-live-label risk is $43/52$, attained by
$Q=\left(\begin{smallmatrix}2/3&1/3\\0&1\end{smallmatrix}\right)$.
For this same angle, two live labels improve on one if and only if $q\alpha>1/3$. This is a phase boundary for the declared random-horizon prediction task, not the unresolved intermediate-memory deterministic-release law.
\end{theorem}
\begin{proof}
At the least budget, Lemma~\ref{lem:boundary-herald} reduces every feasible observer to first-herald commitment. Its two-state failure matrix has the form
$Q=\left(\begin{smallmatrix}1-a&a\\b&1-b\end{smallmatrix}\right)$, $0\le a,b\le1$, from initial label zero. Write $u_i\in[-1,1]$ for the expected decoded sign at label $i$. Its risk is
\[
 2-2v\cdot u,\qquad
 v=p\sum_{n\ge0}x^n\cos(n\vartheta)e_0Q^n.
\]
The series is absolutely convergent because $x<1$. If $a+b=0$, then $v=(A,0)$. Otherwise put $\lambda=1-a-b$ and $z=a/(a+b)$. The elementary two-state decomposition gives
\[
 v_0=(1-z)A+zB,\quad v_1=z(A-B),\qquad
 B=p f_c(x\lambda)>0.
\]
Also $A>0$, so $v_0>0$. Optimizing the two decoder signs therefore gives reward $A$ if $B\le A$, and $A+2z(B-A)$ otherwise. For a fixed $\lambda$, the largest feasible $z$ is one when $\lambda\ge0$ and $1/(1-\lambda)$ when $\lambda<0$.

For $c\ge0$, differentiation gives
$f_c'(t)=(c-2t+ct^2)/(1-2ct+t^2)^2$. Its maximum on $[-x,x]$ occurs at $\min\{x,t_*\}$. When $x\le t_*$, $B\le A$ for every matrix. Otherwise the maximizing $\lambda=t_*/x$ is nonnegative and permits $z=1$. Thus the upper $2\max B-A$ is attainable, and $f_c(t_*)=(1+s)/(2s)$ gives the formula.

For $c=-k<0$, nonnegative $\lambda$ is optimized at zero. For negative $\lambda$ set $t=-x\lambda\in[0,x]$ and choose $z=1/(1+t/x)$. The improvement divided by two is
\[
 \frac{p x}{1+2kx+x^2}
       \frac{x+k-(1+kx)t}{1-2kt+t^2}.
\]
Its derivative has the sign of
$(1+kx)t^2-2(x+k)t+2k(x+k)-(1+kx)$.
The larger root exceeds $x$ because $(x+k)/(1+kx)>x$ for $k>0$. Consequently the maximum on $[0,x]$ is at the smaller root clipped to this interval; that root is the displayed $t$. The endpoint zero includes the best nonnegative-$\lambda$ choice. This proves global optimality over all two-state matrices and all randomized decoders. Substitution gives $A=7/13$, $\max B=9/16$ in the specified algebraic case, hence $2-2(2\cdot9/16-7/13)=43/52$. The strict improvement condition follows from the positive-angle cases. No numerical optimization is used in this proof.
\end{proof}

\subsection{A finite-precision synthesis bound for unary rational channels}
There is also a positive computational conclusion at the boundary, rather than only a nonattainment theorem. This statement concerns one known prepared channel and a fixed finite dictionary; the full robust synthesis problem is not asserted to have the same complexity.

\begin{theorem}[Polynomial prefix synthesis at a bounded herald rate]\label{thm:boundary-algorithm}
Suppose $Y$ has one letter, the prepared density matrix and the channel superoperator in real Hermitian coordinates have rational entries, and the $m$ loss observables have rational entries and lie in $[0,LI]$. Let $p$ and a requested $0<\epsilon<L$ be rational. From these data one can construct an autonomous observer of exact cost $1/p$ and risk at most $V+\epsilon$, using
\[
 T=\max\left\{1,\left\lceil\frac{\log(L/\epsilon)}{-\log(1-p)}\right\rceil\right\}
\]
live labels. The integer $T$ may instead be chosen by exact successive rational tests of $L(1-p)^T\le\epsilon$, without evaluating a logarithm.

After a common-denominator preprocessing, suppose the data entries have denominator dividing $D$ and numerators and $D$ have at most $B$ bits; include the rational herald rate, loss bound and tolerance in the input size. The synthesis uses $O(T(d^4+md^2))$ rational coordinate operations on integers with $O(TB+\log d)$ bits, in addition to the scalar stopping test. A direct implementation has $O(d^2(TB+\log d))$ bits of working storage, apart from its output program, and a transition/readout table of $O(T(\log(T+1)+\log(m+1)))$ bits. Thus, for a herald rate bounded away from zero, this is a polynomial bit-time construction in the rational input size and $\log(1/\epsilon)$, using standard integer arithmetic. It is not a polynomial bound uniform as $p\downarrow0$, a matching lower bound, or a complexity claim for continuum-model robust optimization.
\end{theorem}
\begin{proof}
Compute $\rho_n=\Phi^n(\rho_0)$ for $0\le n<T$. At each step evaluate the $m$ rational numbers $\tr(F_r\rho_n)$ and choose a least one, resolving ties by the fixed dictionary order. Use a chain of $T$ live labels; on failure advance until the last label, which repeats. At a herald commit the decision stored at that label. The rule is optimal on all terminal words of length less than $T$. Its excess loss is at most $L$ on the remaining terminal mass $q^T$, while its call count is still geometric with mean $1/p$. This proves the guarantees.

A channel application is a matrix--vector product in $d^2$ real Hermitian coordinates, requiring $O(d^4)$ scalar operations. Evaluating all decisions costs $O(md^2)$ more. Use real and imaginary matrix-entry coordinates, with off-diagonal entries paired, so all physical state coordinates have absolute value at most one. After $n$ applications their denominator divides $D^{n+1}$. Keeping this common denominator explicitly gives integer numerators of $O((n+1)B)$ bits; a loss inner product adds only $O(B+\log d)$ bits. The same size bound follows directly during multiplication and summation, and comparison of decisions uses integer comparisons at the shared denominator. Store two state vectors and the current best decision; the stated working bound includes the growing state numerators. The input itself is read-only and not included in working storage. The output table stores a successor and a dictionary index per live label, with no stored floating phase or external clock. Standard integer addition and multiplication have polynomial bit cost, giving the final claim. If the exact scalar test is used, the denominator and numerator of $q^T$ also have $O(TB)$ bits. The displayed bound on $T$ is logarithmic in precision only when $p$ has the stated positive lower bound.
\end{proof}

The qubit in Theorem~\ref{thm:qubit-boundary} belongs to this rational-superoperator class: although $U$ has algebraic entries, its conjugation map has rational coefficients with denominator five. Its exact optimum requires unbounded classical memory, but its boundary-optimal risk has efficiently computable finite-memory approximants. The stored number of live-label bits is $\lceil\log_2T\rceil$; this does not replace the separate program and arithmetic storage accounts above.

The algebraic rotating-qubit idea and nonperiodic quantum acceptance events have established antecedents in quantum automata \cite{Mereghetti24}. We do not claim priority for that mechanism. Here it is compiled into a paid, single-command prediction experiment whose least feasible cost is attained, whose full-history bounded-risk optimum is attained, and whose exact risk is nevertheless unattainable with every fixed finite classical memory. The same boundary admits exponentially accurate finite-label approximations. The transition-table account excludes the fixed numerical dictionary and primitive description, whose encoding lengths must also be added when measuring total program size. These facts distinguish optimum attainment, fixed-memory gaps, and density, rather than inferring one from another. They do not settle the sharp intermediate-memory phase-type timing curve of Theorem~\ref{thm:timer-frontier}, the general boundary-density question, or a universal unknown-model predictive quotient.

\section{A timing task with an exact attainment threshold}\label{sec:timing-memory}
The complete stopped law can distinguish controllers whose terminal outputs are identical. A basic instance is a delayed-release protocol: a stateless physical device emits the same productive report at every unit-cost call, and the observer is asked to commit at a prescribed call. There is one legal command, a constant decoder and a constant audit. An autonomous observer has no free clock, call counter or persistent random tape. Its retained label is its only intercall state. Thus finite live-state implementations are discrete phase-type stopping mechanisms. Phase-type timing and its optimal variability have classical antecedents; the continuous-time coefficient-of-variation theorem of Aldous--Shepp \cite{AldousShepp87} is a different objective from the complete-law, discrete-call approximation considered here.

Fix an integer $K\ge2$ and let the target stopped law be commitment exactly at call $K$. For a proper observer with stopping time $\tau\ge1$, its total-variation discrepancy from this target is $1-\Pr(\tau=K)$. Write $\Delta_{N,M}(K)$ for the smallest discrepancy at expected-call budget $N\ge1$ with at most $M$ autonomous live labels. Event stopping may reuse one live label for the constant decoder. State stopping requires an additional disjoint terminal label. Let $\Delta_{N,\infty}(K)$ denote the full-history optimum. This notation counts \emph{live} labels, not total labels, in both conventions.

\begin{theorem}[Delayed release: budget curve, exact state threshold and pumping gap]\label{thm:timer-frontier}
For every $K\ge2$ and $N\ge1$,
\begin{equation}\label{eq:timer-unrestricted}
 \Delta_{N,\infty}(K)=\frac{(K-N)_+}{K-1}.
\end{equation}
For one live label, put $n=\min\{N,K\}$. The exact optimum is
\begin{equation}\label{eq:timer-one-state}
 \Delta_{N,1}(K)=1-\frac1n\left(1-\frac1n\right)^{K-1}.
\end{equation}
At every $N>1$, the least live-label count attaining \eqref{eq:timer-unrestricted} is exactly $K$. Consequently its least total count is $K$ for event stopping and $K+1$ for state stopping. At $N=1$, immediate commitment with one live label already attains the discrepancy one.

More quantitatively, for $M<K$, every feasible autonomous observer with $s=\Pr(\tau=K)$ obeys
\begin{equation}\label{eq:timer-pumping}
 s+\frac{s^2}{M^{2K-2}}\le1,
 \qquad
 (K-1)s+\frac{K s^2}{M^{2K-2}}\le N-1.
\end{equation}
These explicit bounds force a strict gap from the unrestricted optimum at every $N>1$. For $N\ge K$ they give the budget-independent bound
$\Delta_{N,M}(K)\ge 1/(4M^{2K-2})$.
The one-label curve and $K$-label attainment threshold are exact; no sharpness of the intermediate-$M$ quantitative bound is asserted.
\end{theorem}
\begin{proof}
Put $s=\Pr(\tau=K)$. Since $\tau\ge1$,
$\E\tau\ge1+(K-1)s$. Thus $s\le\min\{1,(N-1)/(K-1)\}$. This is attained by stopping at call one with probability $1-s$ and at call $K$ with probability $s$. It is an actual causal construction: at the first call either commit or enter a finite deterministic chain, and then advance one counted label per call until commitment. It uses $K$ live labels when $s>0$, with no external clock or surviving seed. For $N\ge K$ use the deterministic chain. This proves \eqref{eq:timer-unrestricted} and its upper state count. Equivalently, the bounded score $\one_{\{\tau\ne K\}}$ witnesses this complete-law discrepancy in Theorem~\ref{thm:complete-stopped-order}; the constant terminal-output marginal would not witness it.

With one live label the post-report halt probability is one constant $p>0$, so $\tau$ is geometric on the positive integers. The budget says $p\ge1/N$, and its target mass is $p(1-p)^{K-1}$. Its derivative on $(0,1)$ has the sign of $1-Kp$. The maximum over $[1/N,1]$ is therefore at $p=1/\min(N,K)$. The endpoints give the same formula, proving \eqref{eq:timer-one-state}.

For the lower state count and the quantitative bound, let $Q$ be the substochastic transition matrix on the live labels and let $h=(I-Q)\one$ be its halt vector. The initial live label is prescribed. A path that halts at call $K$ contains $K$ live labels before its calls and at most $M^{K-1}$ possible label sequences. If $s>0$, one such path has probability $q\ge s/M^{K-1}$. When $M<K$, its label sequence repeats a label. The intervening live cycle has length $\ell\ge1$ and transition-product probability $r$. Every other factor in the original path is at most one, so $r\ge q$. Traverse that same cycle once more and follow the original remainder to halt. This is a legal path of length $K+\ell$, of probability $qr\ge q^2$. It is disjoint from all length-$K$ paths. Normalization yields $1\ge s+q^2$, proving the first inequality. The identity $\E\tau=1+\E(\tau-1)$ yields
\[
 \E\tau\ge1+(K-1)s+(K+\ell-1)q^2
          \ge1+(K-1)s+Kq^2,
\]
which proves the second. Zero-probability paths need no conditioning, and the case $s=0$ is immediate.

For $1<N<K$, the positive quadratic term makes the largest permitted $s$ strictly smaller than $(N-1)/(K-1)$; for $N\ge K$ the first inequality makes it strictly smaller than one. Thus fewer than $K$ live labels cannot attain the unrestricted curve, even as an infimum. If $s\ge1/2$, the first inequality gives $1-s\ge1/(4M^{2K-2})$; if $s<1/2$, that bound is weaker than $1-s>1/2$. This gives the last assertion. The globally typed state-stopped interface has a separate terminal label, whereas the final event permits reuse in event stopping, yielding the two total counts.
\end{proof}

This temporal lower bound is not a command-domain incompatibility count: every live context has the very same legal command. It comes from repeated autonomous state and the extra observable stopping histories forced by a positive cycle. A clock supplied as a physical input, a nonconstant report encoding elapsed time, or a persistent external seed changes the interface and invalidates that lower. The result is an exact delayed-release consequence with matching budget and attainment counts, not a solution of the entire intermediate-state phase-type optimization or a claim to settle a previously recognized open timing problem.

\section{The causal task-order transfer theorem}\label{sec:central}
Suppose that every terminating state-stopped competitor reaches at least $B_0\ge1$ live labels under $\theta_*$. Such a certificate must concern actual physical command compatibility, not printed action names. Define
\begin{equation}\label{eq:widths}
\begin{array}{c|cc}
 i&m_-(M)&m_+(M)\\\hline
 \cp&M&\lfloor M/D_0\rfloor\\
 \ev&M&\lfloor M/D_0\rfloor\\
 \st&M-B_0&\lfloor(M-B)/D_0\rfloor .
\end{array}
\end{equation}
The event upper also requires $M\ge B$. The trivial state-stopped lower $B_0=1$ is valid for a procedure initially waiting to make a call.

\begin{theorem}[Affine task-order transfer with a paid reference and common construction]\label{thm:transfer}
Assume all the prepared-kernel, full-report, task-order, effort, occupation, resource and common-construction hypotheses of Section~\ref{sec:data}. Fix $N,M$, put $b_-=N/\kappa_-$, $b_+=N/\kappa_+$, and suppose $b_-\ge1$ and $m_-(M)\ge1$. Set
\[
 t_N=(1-A_N)_+/u_-.
\]
If $t_N>0$, every one of the three risks satisfies
\begin{equation}\label{eq:central-lower}
 R_i^\Theta(N,M)\ge
 \bigl[l_-t_N\Q_{m_-(M)}(b_-/t_N)-E_N\bigr]_+.
\end{equation}
If $t_N=0$, the asserted lower is zero, with no evaluation of $\Q$ at infinity. An empty legal class still has risk $+\infty$.

Suppose also $m_+(M)\ge1$, $b_+\ge1$, and, for event stopping, $M\ge B$. If
\begin{equation}\label{eq:beta}
 \beta_N=\frac{l_+b_+}{1+b_+a_+}\ge1,
\end{equation}
one common Borel program with actual expected cost at most $N$ in every model attains
\begin{equation}\label{eq:central-upper}
 R_i^\Theta(N,M)\le
 \frac{u_+(1+b_+a_+)}{l_+}\Q_{m_+(M)}(\beta_N)+b_+e_+.
\end{equation}
It commits almost surely. If $\beta_N<1$, the always-accept first-attempt construction instead gives $R_i^\Theta(N,M)\le u_+\Q_{m_+(M)}(1)+e_+$ at cost at most $\kappa_+\le N$.

The lower is against every legal competing strategy and uses one actual witness. The upper is one parameter-independent strategy for all models. The two effort constants need not agree. Neither direction assumes that the comparison experiment belongs to $\Theta$ or that a normalized predictive update is uniformly contracting.
\end{theorem}

\begin{corollary}[Scale-uniform relative transfer]\label{cor:relative}
Under Theorem~\ref{thm:transfer}, suppose $A_N=E_N=a_+=e_+=0$. If $l_+b_+\ge1$, then
\begin{equation}\label{eq:relative}
 \frac{l_-}{u_-}\Q_{m_-(M)}(u_-b_-)
 \le R_i^\Theta(N,M)
 \le\frac{u_+}{l_+}\Q_{m_+(M)}(l_+b_+).
\end{equation}
The factors $l_\pm,u_\pm$ may be fixed nonunit constants, independent of budget and resolution. In particular no term proportional to $b_\pm$ times a fixed absolute calibration error is required in this relative-order class.
\end{corollary}
\begin{proof}
Substitute $t_N=1/u_-$ and $\beta_N=l_+b_+$ in the theorem.
\end{proof}

\begin{corollary}[Separate exact checkpoint and autonomous laws]\label{cor:exact}
Suppose the certificates are exact: $l_\pm=u_\pm=1$, all defects vanish, $\kappa_-=\kappa_+=\kappa$, and exports are the identity. Then, for $N\ge\kappa$,
\[
 R_\cp^\Theta(N,m)=\Q_m(N/\kappa).
\]
The constructed autonomous upper uses $B+m$ state-stopped labels, or $\max(B,m)$ event-stopped labels. If the mandatory live count is sharp, $B_0=B$, then
\begin{align}
 R_\st^\Theta(N,B+m)&=\Q_m(N/\kappa)\quad(m\ge1),\label{eq:typed-exact}\\
 R_\ev^\Theta(N,M)&=\Q_M(N/\kappa)\quad(M\ge B).\label{eq:event-exact}
\end{align}
The subtraction by $B$ belongs only to the state-stopped model. Nonidentity exports require the expanded dictionary in \eqref{eq:widths}, unless an optimal dictionary is invariant under all exports. These exact stopping distinctions are inherited from the complete companion and are reproved here to fix the resource semantics.
\end{corollary}
\begin{proof}
The same lower and upper profile coincide in each asserted interface. Identity exports give $D_0=1$. The packing and converse widths are proved in Lemma~\ref{lem:roles}. No lower on live states is used to forbid event-stopped reuse.
\end{proof}

The familiar absolute-defect theorem is also a specialization: take $l_\pm=u_\pm=1$ while retaining $A_N,E_N,a_+,e_+$. Thus relative and absolute calibration are composed in one statement. The chosen metric is the common scored task throughout; it is not changed between the converse and the construction.

\section{Proof of the central theorem}
\begin{lemma}[Profile attainment]\label{lem:attainment}
For every finite $m\ge1$ and $b\ge1$, the minimum in \eqref{eq:profile} is attained. The profile is nonincreasing in $m,b$ and left-continuous in $b>1$.
\end{lemma}
\begin{proof}
Write a dictionary as an ordered tuple in the compact set $H^m$, permitting repetitions. The feasible densities are a weak-star closed subset of the radius-$b$ ball of $L^\infty(\nu)$: positivity, the upper bound and the mass equality are each closed under pairing with $L^1(\nu)$. They are therefore weak-star compact. As the tuple varies, its squared-distance function varies uniformly on $H$, since all points belong to a bounded set. Along a convergent product net, split the objective difference into this uniform variation and the weak-star pairing against the fixed bounded measurable loss at the limit tuple. This proves continuity of the objective and attainment. Nets avoid an unwarranted sequential compactness assumption. Monotonicity follows from inclusion of feasible sets. For left continuity take an optimizer $h$ at $b>1$ and let $h_\lambda=(1-\lambda)h+\lambda$. It has unit mass and cap $b_\lambda=(1-\lambda)b+\lambda<b$. Its objective tends to the optimum as $\lambda\downarrow0$, while monotonicity bounds the optimum at $b_\lambda$ between the optimum at $b$ and this objective.
\end{proof}

\begin{lemma}[Actual-history comparison measure]\label{lem:comparison}
Fix any robustly admissible competitor of cost at most $N$, with its fixed terminal dictionary $C$. There are finite nonnegative measures $\xi_c$ on $X$, $c\in C$, such that $\xi=\sum_c\xi_c$ obeys
\begin{align}
 \xi&\le (\E_*A_\tau)\nu\le b_-\nu,\label{eq:measure-cap}\\
 1&\le u_-\xi(X)+A_N,\label{eq:measure-mass}\\
 \E_*\|Z-\widehat r\|^2
 &\ge l_-\sum_{c\in C}\int d_c\,d\xi_c-E_N.
 \label{eq:measure-risk}
\end{align}
The measure $\xi$ is not asserted to be the terminal law of any comparison execution.
\end{lemma}
\begin{proof}
Adjoin an independent uniform random variable to the current report on both sides. This is current-call randomness, not a seed retained between calls; it is sufficient to represent the competitor's randomized output and stopping decision. A finite-output kernel has such a representation by cumulative probabilities. At a fixed reached productive history $h$ and chosen command $a$, the post-report decision then gives an acceptance indicator $q_h(y)$ and an output $r_h(y)\in C$ on acceptance. Give $r_h$ an arbitrary fixed value on rejection. The notation $y$ now includes the current independent random variable. Integrating the original certificates over that variable preserves their constants. Dependence on the entire earlier history is permitted in these functions.

For each Borel $D\subset X$, set
\begin{equation}\label{eq:xi-definition}
 \xi_c(D)=\E_*\sum_{j\le\tau:\ j\ {
m productive}}
 \int\one_{\{\pi_{H_{j-1},a_j}(y)\in D\}}
 q_{H_{j-1}}(y)\one_{\{r_{H_{j-1}}(y)=c\}}
 P^0_{H_{j-1},a_j}(dy,H).
\end{equation}
The outer expectation is the actual pre-call occupation at $\theta_*$. All terms are jointly measurable and nonnegative. Tonelli defines the measures, whose total is finite because $\E_*A_\tau\le b_-$. Summing over $c$ and dropping the acceptance indicator gives \eqref{eq:measure-cap}: each comparison row has mark law $\nu$.

The analogous sum with the actual rows instead of $P^0$ has mass exactly one, because the actual procedure commits exactly once almost surely. Apply \eqref{eq:lower-mass} at each actual reached history, then sum. The occupation certificate yields \eqref{eq:measure-mass}. Apply \eqref{eq:lower-task} in the same way to the disjoint actual commitment events. Their expected losses sum to the actual risk. The comparison contribution is $d_{r_h(y)}(\pi_{h,a}y)$ by \eqref{eq:full-moments}; this identity conditions on the full current report and therefore does not overlook extra information seen by the competitor. The accumulated loss defect is at most $E_N$, giving \eqref{eq:measure-risk}. One may first truncate the number of calls. Monotone convergence for the nonnegative contributions and the finite defect bounds justify the limit.

No actual past is replaced by a comparison past. In particular the competitor need not terminate, or obey a cost bound, if hypothetically run under the comparison experiment. This distinction is why the argument applies when that experiment is not an operating model.
\end{proof}

\begin{remark}[Why the full report is required]
A constant mark can conceal a visible fair bit which also equals the future audit. A comparison using only that mark has $f=0,w=1$, but a two-value decoder reading the bit predicts it with zero loss. The hypotheses therefore condition on the full report seen by the competitor; matching a compressed mark marginal alone cannot establish the lower.
\end{remark}

\begin{lemma}[Live roles, output roles and physical compatibility]\label{lem:roles}
A program with $\ell$ live roles and $t$ output roles has a state-stopped implementation with $\ell+t$ labels and an event-stopped implementation with $\max(\ell,t)$ labels, preserving every path cost and complete stopped scored law. Conversely an event-stopped $M$-label program has at most $M$ decoder values. A state-stopped $M$-label program with at least $B_0$ reached live labels has at most $M-B_0$ terminal decoder values.

For a physical live-state certificate, consider finitely many mandatory pre-call contexts $I$ with no free context input. Suppose the commands legal at a history of context $i$ are exactly $\mathcal U_i$, and every terminating competitor visits every context with positive probability under the same witness. This positive probability may depend on the competitor; no uniform lower bound over competitors is required. Let $b_{\rm phys}$ be the least size of a cover of $I$ by subsets $J$ with $\bigcap_{i\in J}\mathcal U_i\ne\varnothing$. Then at least $b_{\rm phys}$ live labels are required. This certificate is invariant under duplicate names and bijective stateless recodings of the same physical commands.
\end{lemma}
\begin{proof}
For state stopping take disjoint copies of the live and terminal role sets and type them zero and one. For event stopping embed the live and terminal role sets separately into one alphabet of size $\max(\ell,t)$. On continuing transitions use the live embedding; on commitment emit the final event and use the terminal embedding. A command row is read only on a continuing branch, while the decoder is read only after commitment. Sharing their label therefore causes no conflict and uses no continuing mode bit. Couple the original and recoded coins. Every finite stopped prefix and every path cost agree, and almost-sure termination gives equality of the stopped scored laws. A fixed decoder has one value per retained label. In state stopping the global halt map excludes each reached live label from the output set, proving its subtraction. Event stopping has no such exclusion.

A live label reached in contexts $J$ must have one command distribution supported on all the corresponding $\mathcal U_i$. Since $J$ is finite, the union of the exceptional null command sets is null; hence their intersection supports a probability and is nonempty. The sets of contexts served by the labels cover $I$, proving the bound. Conversely a nonempty intersection contains a command and hence a Dirac command distribution. Replacing names without changing executed kernels, legal domains, costs and reports does not change these intersections. A translator which itself remembers a hidden phase is not such a recoding; its memory must be charged.
\end{proof}

\begin{proof}[Proof of Theorem~\ref{thm:transfer}]
Fix a competing program. Lemma~\ref{lem:roles} bounds its decoder width by $m_-(M)$. Let $s=\xi(X)$ in Lemma~\ref{lem:comparison}. If $t_N>0$, then \eqref{eq:measure-mass} gives $s\ge t_N$. Scale the measure to $\xi'=(t_N/s)\xi$. It has mass $t_N$ and remains bounded above by $b_-\nu$. Thus $h=(d\xi'/d\nu)/t_N$ has unit mass and cap $b_-/t_N$. Nearest-dictionary projection and nonnegativity imply
\[
 \sum_{c\in C}\int d_c\,d\xi_c
 \ge\int [w+\dist(f,C)^2]\,d\xi
 \ge t_N\Q_{m_-(M)}(b_-/t_N).
\]
Together with \eqref{eq:measure-risk} this proves the asserted lower for this competitor. Its worst-model risk dominates its risk at $\theta_*$. Taking the infimum over competitors gives \eqref{eq:central-lower}, without exchanging an infimum and a supremum. Nonnegative loss gives the zero truncation and the stated case $t_N=0$.

For the upper assume \eqref{eq:beta}. Choose an attained optimizer $(C,h)$ at width $m_+(M)$ and cap $\beta_N$. Let $r_C(x)$ be the first nearest point to $f(x)$ in a fixed ordering of $C$. After each attempt of $\mathcal G$, accept $x$ with probability $h(x)/\beta_N$, export $T_Jr_C(x)$ and commit; otherwise return to the start routing label. All these operations are Borel. The choice $J$ is used within the current attempt; any part retained thereafter is charged through the live role or exported decoder, not a free seed. The exports are assumed jointly measurable and take at most $D_0|C|$ distinct values, so their finite decoder and transitions are Borel as well.

Conditional on every surviving attempt history, the acceptance probability is at least
\begin{equation}\label{eq:hazard}
 p=\frac{l_+}{\beta_N}-a_+=\frac1{b_+}>0.
\end{equation}
Consequently the probability of reaching attempt $j+1$ is at most $(1-p)^j$, by induction on conditional survival probabilities. In particular commitment is almost sure and the expected number of attempts is at most $1/p=b_+$. Each begun attempt has conditional expected cost at most $\kappa_+$, not merely a cost conditional on a favorable outcome. Summing the predictable attempt costs gives actual expected cost at most $\kappa_+b_+=N$.

The conditional unnormalized loss contributed by an attempt is at most $u_+\Q_{m_+(M)}(\beta_N)/\beta_N+e_+$. Sum its contributions on the disjoint terminal events, using the same survival bound. The result is at most
\[
 b_+\left[\frac{u_+}{\beta_N}\Q_{m_+(M)}(\beta_N)+e_+\right],
\]
which is \eqref{eq:central-upper}. The program has $B$ live roles and at most $D_0m_+(M)$ terminal roles. Lemma~\ref{lem:roles} gives exactly the sufficient widths in \eqref{eq:widths}. A checkpoint counts only the terminal width. No physical renewal law was used: the conditional hazard controls survival even if the hidden state after rejection differs arbitrarily from its former state.

If $\beta_N<1$, commit after the first attempt, and use an optimizer at cap one. The loss certificate with $q=1$ gives $u_+\Q_{m_+(M)}(1)+e_+$, and the whole-attempt cost is at most $\kappa_+$. This proves all remaining assertions.
\end{proof}

\section{Scope and continuing mathematical problems}\label{sec:r66-scope}
The closed performed-law family and its decision tests concern a common finite-branching interface, a positive conditional call-cost floor, bounded terminal loss and a fixed actual dictionary. They do not imply uniform expected cost-tail integrability or finite classical-memory attainment at the least budget. The heralded-channel results classify that latter question within their prepared-channel interface and prove both positive finite-memory realizations and a quantum obstruction. The prefix algorithm has its own unary rational-data assumptions. No result here turns arbitrary quantum cutpoint acceptance into exact deterministic prediction.

The geometric route remains independent: a commandwise lower witness on actual histories and one common paid upper construction yield the comparison profile, not conversely for every primitive experiment. Its width, rank, singular-range, regular, mixed and atomic consequences are proved in the current supplement. The separately retained nonintervening-sensing, ordered-measurement, reset-communication and rank-constrained-discrimination lines are not consequences inferred from the heralded example. The predictive quotient still has its stated fixed-law scope.

Open questions include an intrinsic general criterion for geometric exactness, common unknown-model predictive realizations, boundary finite-memory density beyond the class treated here, the sharp intermediate-memory phase-type timing frontier, and matching program, precision, workspace and time lower bounds. The new weighted and tied criterion requires a single command, a uniformly proper first-herald instrument, a fixed dictionary and the stated cost-support coverage. It does not classify arbitrary controlled boundary stopping. The positive realization theorem classifies exact stochastic implementation on this optimal face; the finite algebraic minimal-synthesis result has a supplied finite linear response and a proved size upper. It is not a decision procedure for every unbounded positive-realization problem. The exact three-label calibration threshold does not solve the older general-$M$ squared-sign or intermediate phase-type timing frontiers. The original title and research objective are retained without claiming that these distinct open questions have been settled.
\begin{thebibliography}{99}
\bibitem{Blackwell} D. Blackwell, \emph{Equivalent comparisons of experiments}, Ann. Math. Statist. \textbf{24} (1953), 265--272. DOI: 10.1214/aoms/1177729032.
\bibitem{CGM97} J. A. Cuesta-Albertos, A. Gordaliza and C. Matr\'an, \emph{Trimmed $k$-means: an attempt to robustify quantizers}, Ann. Statist. \textbf{25} (1997), 553--576. DOI: 10.1214/aos/1031833664.
\bibitem{GG99} L. A. Garc\'ia-Escudero and A. Gordaliza, \emph{Robustness properties of $k$ means and trimmed $k$ means}, J. Amer. Statist. Assoc. \textbf{94} (1999), 956--969. DOI: 10.1080/01621459.1999.10474200.
\bibitem{SC} C. R. Shalizi and J. P. Crutchfield, \emph{Computational mechanics: pattern and prediction, structure and simplicity}, J. Statist. Phys. \textbf{104} (2001), 817--879. DOI: 10.1023/A:1010388907793.
\bibitem{BC} N. Barnett and J. P. Crutchfield, \emph{Computational mechanics of input--output processes: structured transformations and the $\varepsilon$-transducer}, J. Statist. Phys. \textbf{161} (2015), 404--451; arXiv:1412.2690.
\bibitem{AMTW} A. Ambainis, M. Mosca, A. Tapp and R. de Wolf, \emph{Private quantum channels}, Proc. 41st Annual Symposium on Foundations of Computer Science (2000), 547--553. DOI: 10.1109/SFCS.2000.892142.
\bibitem{Buscemi} F. Buscemi, \emph{Comparison of quantum statistical models: equivalent conditions for sufficiency}, Comm. Math. Phys. \textbf{310} (2012), 625--647. DOI: 10.1007/s00220-012-1421-3.
\bibitem{BR14} N. Bertschinger and J. Rauh, \emph{The Blackwell relation defines no lattice}, Proc. IEEE International Symposium on Information Theory (2014), 2479--2483. DOI: 10.1109/ISIT.2014.6875280; arXiv:1401.3146.
\bibitem{Banerjee05} A. Banerjee, X. Guo and H. Wang, \emph{On the optimality of conditional expectation as a Bregman predictor}, IEEE Trans. Inform. Theory \textbf{51} (2005), 2664--2669. DOI: 10.1109/TIT.2005.850145.
\bibitem{Stengel96} B. von Stengel, \emph{Efficient computation of behavior strategies}, Games Econom. Behav. \textbf{14} (1996), 220--246. DOI: 10.1006/game.1996.0050.
\bibitem{Balder89} E. J. Balder, \emph{On compactness of the space of policies in stochastic dynamic programming}, Stochastic Process. Appl. \textbf{32} (1989), 141--150. DOI: 10.1016/0304-4149(89)90058-6.
\bibitem{Barry14} J. Barry, D. T. Barry and S. Aaronson, \emph{Quantum partially observable Markov decision processes}, Phys. Rev. A \textbf{90} (2014), 032311. DOI: 10.1103/PhysRevA.90.032311; arXiv:1406.2858.
\bibitem{Ying18} S. Ying and M. Ying, \emph{Reachability analysis of quantum Markov decision processes}, Inform. Comput. \textbf{263} (2018), 31--51. DOI: 10.1016/j.ic.2018.09.001; arXiv:1406.6146.
\bibitem{Saldi25} N. Saldi, S. Sanjari and S. Y\"uksel, \emph{Quantum Markov decision processes: general theory, approximations, and classes of policies}, SIAM J. Control Optim. \textbf{63} (2025), 1156--1183. DOI: 10.1137/24M1643864; arXiv:2402.14649.
\bibitem{Seidenberg54} A. Seidenberg, \emph{A new decision method for elementary algebra}, Ann. of Math. (2) \textbf{60} (1954), 365--374. DOI: 10.2307/1969640.
\bibitem{Chatterjee15} K. Chatterjee, M. Chmel\'ik, R. Gupta and A. Kanodia, \emph{Optimal cost almost-sure reachability in POMDPs}, Proc. AAAI Conference on Artificial Intelligence \textbf{29} (2015). DOI: 10.1609/aaai.v29i1.9683.
\bibitem{Andriushchenko23} R. Andriushchenko, A. Bork, M. \v{C}e\v{s}ka, S. Junges, J.-P. Katoen and M. Mac\'ak, \emph{Search and explore: symbiotic policy synthesis in POMDPs}, Computer Aided Verification (2023); arXiv:2305.14149v2.
\bibitem{Halmos69} P. R. Halmos, \emph{Two subspaces}, Trans. Amer. Math. Soc. \textbf{144} (1969), 381--389. DOI: 10.1090/S0002-9947-1969-0251519-5.
\bibitem{Kessler21} D. A. Kessler, E. Barkai and K. Ziegler, \emph{First-detection time of a quantum state under random probing}, Phys. Rev. A \textbf{103} (2021), 022222. DOI: 10.1103/PhysRevA.103.022222; arXiv:2012.01763.
\bibitem{DelVecchio26} G. Del Vecchio Del Vecchio and S. N. Majumdar, \emph{Optimal detection of quantum states via projective measurements}, J. Phys. A: Math. Theor. \textbf{59} (2026), 035001. DOI: 10.1088/1751-8121/ae34fe; arXiv:2509.08556.
\bibitem{Companion} Q. Qi, \emph{General Theta Foundations I: Acquired Geometry and Causal Resource Transfer}, complete technical companion on absolute calibration, exact stopping resources and nonintervening sensing, 2026. The complete retained R58 native source and its independently built companion accompany this manuscript; their identities are fixed in the source-provenance and preservation records.
\bibitem{Tzeng92} W.-G. Tzeng, \emph{A polynomial-time algorithm for the equivalence of probabilistic automata}, SIAM J. Comput. \textbf{21} (1992), 216--227. DOI: 10.1137/0221017.
\bibitem{LiYuYing12} Y. Li, N. Yu and M. Ying, \emph{Termination of nondeterministic quantum programs}, arXiv:1201.0891 (2012).
\bibitem{Fu24} J. Fu, H. Jiang, M. Xu, Y. Deng and Z.-B. Li, \emph{Algorithmic analysis of termination problems for nondeterministic quantum programs}, arXiv:2402.15827 (2024).

\bibitem{LiuVerdu18} J. Liu and S. Verd\'u, \emph{Rejection sampling and noncausal sampling under moment constraints}, Proc. IEEE International Symposium on Information Theory (2018), 1565--1569. DOI: 10.1109/ISIT.2018.8437857.
\bibitem{GocFlamich24} D. Goc and G. Flamich, \emph{On channel simulation with causal rejection samplers}, Proc. IEEE International Symposium on Information Theory (2024). DOI: 10.1109/ISIT57864.2024.10619339; examined text arXiv:2401.16579v2 (3 May 2024).
\bibitem{LiLiu26} Z. Li and T. Liu, \emph{Ranking experiments under sequential sampling}, arXiv:2608.19897v2 (28 August 2026). Version 1 appeared under the title \emph{The order of binary experiments under endogenous stopping}, with Z. Li as sole author; the version-2 statement is the one compared here.
\bibitem{Ying14pre} S. Ying and M. Ying, \emph{Reachability analysis of quantum Markov decision processes}, arXiv:1406.6146 (2014), Theorem 3.6; published journal version listed above as \cite{Ying18}.

\bibitem{KaraYuksel20} A. D. Kara and S. Y\"uksel, \emph{Robustness to incorrect system models in stochastic control}, SIAM J. Control Optim. \textbf{58} (2020), 1144--1182. DOI: 10.1137/18M1208058; examined text arXiv:1803.06046v7.
\bibitem{Demirci25} Y. E. Demirci, A. D. Kara and S. Y\"uksel, \emph{Sensitivity of filter kernels and robustness bounds to transition and measurement kernel perturbations in partially observable stochastic control}, arXiv:2508.10658v3 (2025).
\bibitem{Watrous18} J. Watrous, \emph{The Theory of Quantum Information}, Cambridge University Press, 2018, Chapter 3, Corollary 3.40.
\bibitem{KRN65} K. Kuratowski and C. Ryll-Nardzewski, \emph{A general theorem on selectors}, Bull. Acad. Polon. Sci. S\'er. Sci. Math. Astronom. Phys. \textbf{13} (1965), 397--403.

\bibitem{Charnes63} A. Charnes, W. W. Cooper and K. O. Kortanek, \emph{Duality in semi-infinite programs and some works of Haar and Carath\'eodory}, Management Sci. \textbf{9} (1963), 209--228. DOI: 10.1287/mnsc.9.2.209.
\bibitem{Shapiro09} A. Shapiro, \emph{Semi-infinite programming, duality, discretization and optimality conditions}, Optimization \textbf{58} (2009), 133--161. DOI: 10.1080/02331930902730070; author preprint, Optimization Online (2008).
\bibitem{Dufour25} F. Dufour and T. Prieto-Rumeau, \emph{Absorbing Markov decision processes and their occupation measures}, SIAM J. Control Optim. \textbf{63} (2025), 676--698. DOI: 10.1137/24M1661959. Related author preprint examined: \emph{A note on weak compactness of occupation measures for an absorbing Markov decision process}, arXiv:2402.10672v2 (2024).
\bibitem{Franc23} V. Franc, D. Prusa and V. Voracek, \emph{Optimal strategies for reject option classifiers}, J. Mach. Learn. Res. \textbf{24} (2023), no. 11, 1--49.
\bibitem{Jencova16} A. Jen\v{c}ov\'a, \emph{Comparison of quantum channels and statistical experiments}, arXiv:1512.07016v2 (2016). The separately indexed shorter version arXiv:1601.06370 was withdrawn in favor of this extended text.
\bibitem{LiElGamal18} C. T. Li and A. El Gamal, \emph{Strong functional representation lemma and applications to coding theorems}, IEEE Trans. Inform. Theory \textbf{64} (2018), 6967--6978; arXiv:1701.02827v4.
\bibitem{AldousShepp87} D. Aldous and L. Shepp, \emph{The least variable phase type distribution is Erlang}, Comm. Statist. Stochastic Models \textbf{3} (1987), 467--473. DOI: 10.1080/15326348708807067.
\bibitem{Sion58} M. Sion, \emph{On general minimax theorems}, Pacific J. Math. \textbf{8} (1958), 171--176. DOI: 10.2140/pjm.1958.8.171.
\bibitem{PiunovskiyZhang24} A. Piunovskiy and Y. Zhang, \emph{Extreme occupation measures in Markov decision processes with an absorbing state}, SIAM J. Control Optim. \textbf{62} (2024), 65--90. DOI: 10.1137/23M1572398. Author text examined: \emph{Extreme occupation measures in Markov decision processes with a cemetery}, arXiv:2307.03158v1.
\bibitem{DufourMixtures} F. Dufour and T. Prieto-Rumeau, \emph{Absorbing Markov decision processes: geometric properties and sufficiency of finite mixtures of deterministic policies}, arXiv:2512.17511; examined HTML at version-1 URL, Theorems 4.3--4.4 and Proposition 4.5.
\bibitem{Mereghetti24} C. Mereghetti, B. Palano and P. Raucci, \emph{Unary quantum finite state automata with control language}, Appl. Sci. \textbf{14} (2024), no. 4, 1490. DOI: 10.3390/app14041490.
\bibitem{RabinScott59} M. O. Rabin and D. Scott, \emph{Finite automata and their decision problems}, IBM J. Res. Develop. \textbf{3} (1959), 114--125. DOI: 10.1147/rd.32.0114.
\bibitem{Jaeger00} H. Jaeger, \emph{Observable operator models for discrete stochastic time series}, Neural Comput. \textbf{12} (2000), 1371--1398. DOI: 10.1162/089976600300015411.
\bibitem{ThonJaeger15} M. Thon and H. Jaeger, \emph{Links between multiplicity automata, observable operator models and predictive state representations---a unified learning framework}, J. Mach. Learn. Res. \textbf{16} (2015), 103--147.
\bibitem{MonrasWinter14} A. Monras and A. Winter, \emph{Quantum learning of classical stochastic processes: the completely-positive realization problem}, LIPIcs \textbf{27}, TQC 2014 (2014), 99--109. DOI: 10.4230/LIPIcs.TQC.2014.99.
\end{thebibliography}

\end{document}
